% !TEX encoding = IsoLatin2
% notwendige Zeile für Mac-Benutzer (muss als Kommentar stehen);
% Windows-Benutzer können die Zeile löschen.

% LaTeX-Vorlage Version 3.1,  Juli 2011
% erstellt von Dr. Andreas Drauschke (andreas.drauschke@technikum-wien.at) und
% Dr. Susanne Teschl (susanne.teschl@technikum-wien.at)
% geringfügig adaptiert von Harald Stockinger
% (harald.stockinger@technikum-wien.at)

\documentclass[a4paper,bibtotoc,oneside]{scrbook}
% Für kurze Arbeiten wäre auch die Dokumentklasse "scrartcl" ausreichend. In
% diesem Fall ist "section" die höchste Ebene ("chapter" gibt es dann nicht).
% \documentclass[a4paper,bibtotoc,oneside]{scrartcl}

% float
\usepackage{float}

% listings
\usepackage{listings}

\lstdefinelanguage{HTML5} {
  language=html,
  otherkeywords={
    % General
    >,
    %
    <!DOCTYPE,
    <a, </a,
    <abbr, </abbr,
    <address, </address,
    <area,
    <article, </article,
    <aside, </aside,
    <audio, </audio,
    <bdi, </bdi,
    <bdo, </bdo,
    <blockquote, </blockquote,
    <body, </body,
    <br,
    <button, </button,
    <canvas, </canvas,
    <caption, </caption,
    <cite, </cite,
    <code, </code,
    <col,
    <colgroup, </colgroup,
    <command, </command,
    <datalist, </datalist,
    <dd, </dd,
    <del, </del,
    <details, </details,
    <dfn, </dfn,
    <dialog, </dialog,
    <div, </div,
    <dl, </dl,
    <dt, </dt,
    <em, </em,
    <embed,
    <fieldset, </fieldset,
    <figcaption, </figcaption,
    <figure, </figure,
    <footer, </footer,
    <form, </form,
    <h1, </h1,
    <h2, </h2,
    <h3, </h3,
    <h4, </h4,
    <h5, </h5,
    <h6, </h6,
    <head, </head,
    <header, </header,
    <hr,
    <html, </html,
    <i, </i,
    <iframe, </iframe,
    <img,
    <input,
    <ins, </ins,
    <kbd, </kbd,
    <keygen,
    <label, </label,
    <legend, </legend,
    <li, </li,
    <link,
    <map, </map,
    <mark, </mark,
    <menu,
    <meta,
    <meter, </meter,
    <nav, </nav,
    <noscript, </noscript,
    <object, </object,
    <ol, </ol,
    <optgroup, </optgroup,
    <option, </option,
    <output, </output,
    <p, </p,
    <param,
    <pre, </pre,
    <progress, </progress,
    <q, </q,
    <rp, </rp,
    <rt, </rt,
    <ruby, </ruby,
    <s, </s,
    <samp, </samp,
    <script, </script,
    <section, </section,
    <select, </select,
    <small, </small,
    <source,
    <span, </span,
    <strong, </strong,
    <style, </style,
    <sub, </sub,
    <summary, </summary,
    <sup, </sup,
    <table, </table,
    <tbody, </tbody,
    <td, </td,
    <textarea, </textarea,
    <tfoot, </tfoot,
    <th, </th,
    <thead, </thead,
    <time, </time,
    <title, </title,
    <tr, </tr,
    <track,
    <u,</u,
    <ul, </ul,
    <var, </var,
    <video, </video,
    <wbr
  },
  sensitive=true,
  alsoletter={<>=-+},
  morecomment=[s]{<!-}{-->},
  tag=[s],
  morestring=[d]'
}

\lstdefinelanguage{JavaScript} {
  morekeywords={
    break,const,continue,delete,do,while,export,for,in,function,
    if,else,import,in,instanceOf,label,let,new,return,switch,this,
    throw,try,catch,typeof,var,void,with,yield
  },
  sensitive=false,
  morecomment=[l]{//},
  morecomment=[s]{/*}{*/},
  morestring=[b]",
  morestring=[d]'
}

\lstdefinelanguage{CSS} {
  morekeywords={
    accelerator,azimuth,background,background-attachment,
    background-color,background-image,background-position,
    background-position-x,background-position-y,background-repeat,
    behavior,border,border-bottom,border-bottom-color,
    border-bottom-style,border-bottom-width,border-collapse,
    border-color,border-left,border-left-color,border-left-style,
    border-left-width,border-right,border-right-color,
    border-right-style,border-right-width,border-spacing,
    border-style,border-top,border-top-color,border-top-style,
    border-top-width,border-width,bottom,caption-side,clear,
    clip,color,content,counter-increment,counter-reset,cue,
    cue-after,cue-before,cursor,direction,display,elevation,
    empty-cells,filter,float,font,font-family,font-size,
    font-size-adjust,font-stretch,font-style,font-variant,
    font-weight,height,ime-mode,include-source,
    layer-background-color,layer-background-image,layout-flow,
    layout-grid,layout-grid-char,layout-grid-char-spacing,
    layout-grid-line,layout-grid-mode,layout-grid-type,left,
    letter-spacing,line-break,line-height,list-style,
    list-style-image,list-style-position,list-style-type,margin,
    margin-bottom,margin-left,margin-right,margin-top,
    marker-offset,marks,max-height,max-width,min-height,
    min-width,-moz-binding,-moz-border-radius,
    -moz-border-radius-topleft,-moz-border-radius-topright,
    -moz-border-radius-bottomright,-moz-border-radius-bottomleft,
    -moz-border-top-colors,-moz-border-right-colors,
    -moz-border-bottom-colors,-moz-border-left-colors,-moz-opacity,
    -moz-outline,-moz-outline-color,-moz-outline-style,
    -moz-outline-width,-moz-user-focus,-moz-user-input,
    -moz-user-modify,-moz-user-select,orphans,outline,
    outline-color,outline-style,outline-width,overflow,
    overflow-X,overflow-Y,padding,padding-bottom,padding-left,
    padding-right,padding-top,page,page-break-after,
    page-break-before,page-break-inside,pause,pause-after,
    pause-before,pitch,pitch-range,play-during,position,quotes,
    -replace,richness,right,ruby-align,ruby-overhang,
    ruby-position,-set-link-source,size,speak,speak-header,
    speak-numeral,speak-punctuation,speech-rate,stress,
    scrollbar-arrow-color,scrollbar-base-color,
    scrollbar-dark-shadow-color,scrollbar-face-color,
    scrollbar-highlight-color,scrollbar-shadow-color,
    scrollbar-3d-light-color,scrollbar-track-color,table-layout,
    text-align,text-align-last,text-decoration,text-indent,
    text-justify,text-overflow,text-shadow,text-transform,
    text-autospace,text-kashida-space,text-underline-position,top,
    unicode-bidi,-use-link-source,vertical-align,visibility,
    voice-family,volume,white-space,widows,width,word-break,
    word-spacing,word-wrap,writing-mode,z-index,zoom
  },
  sensitive=false,
  morecomment=[l]{//},
  morecomment=[s]{/*}{*/},
  morestring=[b]",
  morestring=[d]'
}

\lstdefinestyle{HTML5} {
  language=HTML5,                               % Language of Code
  basicstyle=\footnotesize\ttfamily,                 % Standardfont for Code
  backgroundcolor=\color{White},
  showstringspaces=false,                            % Underline spaces in St.
  keywordstyle=\ttfamily\bfseries\color{Red},  % Keyword Style
  identifierstyle=\ttfamily,                         % Identifier Style
  stringstyle=\ttfamily\color{OliveGreen},           % String Style
  commentstyle=\color{Gray},                      % Comment Style
  xleftmargin=17pt,                                  % Left Margin
  aboveskip=\bigskipamount,
  belowskip=\bigskipamount,
  numbers = left,
  numbersep = 10pt,
  numberstyle = \color{Gray},
  frame = l,
  rulecolor=\color{Gray},
  tabsize=2                                          % Tabsize in Code
}

\lstdefinestyle{JavaScript} {
  language=JavaScript,                               % Language of Code
  basicstyle=\footnotesize\ttfamily,                 % Standardfont for Code
  backgroundcolor=\color{White},
  showstringspaces=false,                            % Underline spaces in St.
  keywordstyle=\ttfamily\bfseries\color{CadetBlue},  % Keyword Style
  identifierstyle=\ttfamily,                         % Identifier Style
  stringstyle=\ttfamily\color{OliveGreen},           % String Style
  commentstyle=\color{Gray},                      % Comment Style
  xleftmargin=17pt,                                  % Left Margin
  aboveskip=\bigskipamount,
  belowskip=\bigskipamount,
  numbers = left,
  numbersep = 10pt,
  numberstyle = \color{Gray},
  frame = l,
  rulecolor=\color{Gray},
  tabsize=2                                          % Tabsize in Code
}

\lstdefinestyle{CSS} {
  language=CSS,                               % Language of Code
  basicstyle=\footnotesize\ttfamily,                 % Standardfont for Code
  backgroundcolor=\color{White},
  showstringspaces=false,                            % Underline spaces in St.
  keywordstyle=\ttfamily\bfseries\color{Red},  % Keyword Style
  identifierstyle=\ttfamily,                         % Identifier Style
  stringstyle=\ttfamily\color{Green},           % String Style
  commentstyle=\color{Gray},                      % Comment Style
  xleftmargin=17pt,                                  % Left Margin
  aboveskip=\bigskipamount,
  belowskip=\bigskipamount,
  numbers = left,
  numbersep = 10pt,
  numberstyle = \color{Gray},
  frame = l,
  rulecolor=\color{Gray},
  tabsize=2                                          % Tabsize in Code
}

% captions
\usepackage{caption}
\DeclareCaptionFont{white}{\color{white}}
\DeclareCaptionFormat{listing}{\colorbox{gray}{\parbox{\textwidth}{#1#2#3}}}
\captionsetup[lstlisting]{format=listing,labelfont=white,textfont=white}

% verlinkte Querverweise im pdf
\usepackage{hyperref}

% deutsche Anpassungen
\usepackage[utf8]{inputenc}
\usepackage[T1]{fontenc}
\usepackage[ngerman]{babel}

% mathematische Symbole
\usepackage{amsmath,amssymb,amsfonts,amstext}

% Kopfzeilen frei gestaltbar
\usepackage{fancyhdr}
\lfoot[\fancyplain{}{}]{\fancyplain{}{}}
\rfoot[\fancyplain{}{}]{\fancyplain{}{}}
\cfoot[\fancyplain{}{\footnotesize\thepage}]{\fancyplain{}{\footnotesize\thepage}}
\lhead[\fancyplain{}{\footnotesize\nouppercase\leftmark}]{\fancyplain{}{}}
\chead{}
\rhead[\fancyplain{}{}]{\fancyplain{}{\footnotesize\nouppercase\sc\leftmark}}

% Farben im Dokument möglich
\usepackage[usenames,dvipsnames]{color}

% enumerate styling
\usepackage{enumitem}

% Schriftart Helvetica
\usepackage{helvet}
\renewcommand{\familydefault}{cmss}

% Graphiken einbinden: hier für pdflatex
\usepackage[pdftex]{graphicx}

\usepackage{array}

% Höhe und Breite des Textkörpers etwas grösser definieren
\setlength{\textheight}{225mm}
\setlength{\textwidth}{1.05\textwidth}

% weniger Warnungen wegen überfüllter Boxen
\tolerance = 9999
\sloppy

% Anpassung einiger Überschriften
\renewcommand\figurename{Abbildung}
\renewcommand\tablename{Tabelle}

\begin{document}

% Kopf- und Fusszeilen initiieren
\pagestyle{fancy}

% Deckblatt:
\thispagestyle{empty}
\begin{picture}(0,0)
\color{white}\sffamily
\put(-101,-749){\includegraphics[width=1.002\paperwidth,
height=\paperheight]{BM_2011.pdf}}
\put(220,-670){\includegraphics[width=0.5\textwidth]{FHTW_Logo_4c.pdf}}
\put(-30, -20){\bfseries\huge MASTER THESIS}
\put(-30,-50){\Large zur Erlangung des akademischen Grades}
\put(-30,-70){\Large \glqq Master of Science in Engineering\grqq}
% Titel des Studienganges einfügen:
\put(-30,-90){\Large im Studiengang Multimedia und Softwareentwicklung}
% Titel der Arbeit einfügen:
% Die Minipage wird gesetzt, damit auch mehrzeilige Titel möglich werden.
\put(-32,-180){
\begin{minipage}{14cm}
  \bfseries\huge HTML5 Canvas vs. Adobe Flash
\end{minipage}
}
% Name der Autorin/des Autors eingeben:
\put(-30,-270){\large Ausgeführt von: Peter Kerschner, BSc.}
% Personenkennzeichen der Autorin/des Autors eingeben:
\put(-30,-290){\large Personenkennzeichen: 1210299029}
% Name der Begutachterinnen/der Begutachter eingeben:
\put(-30,-330){\large 1. BegutachterIn: Dipl.-Ing. (FH) Arthur Michael Zaczek}
\put(-30,-350){\large 2. BegutachterIn: Dipl.-Ing. Mag. Dr. Michael Tesar}
\put(-30,-390){\large Wien, \today} % das Datum des letzten Kompilierens wird
% automatisch eingesetzt
\color{black}
\end{picture}

\newpage


\section*{Eidesstattliche Erklärung}\thispagestyle{empty}
\glqq Ich erkläre hiermit an Eides statt, dass ich die vorliegende Arbeit
selbständig angefertigt habe.
Die aus fremden Quellen direkt oder indirekt übernommenen Gedanken sind als
solche kenntlich gemacht.
Die Arbeit wurde bisher weder in gleicher noch in ähnlicher Form einer anderen
Prüfungsbehörde vorgelegt
und auch noch nicht veröffentlicht. Ich versichere, dass die abgegebene Version
jener im Uploadtool entspricht.\grqq\\[5\baselineskip]
\rule{5cm}{0.2pt}\hfill\rule{5cm}{0.2pt}\\
\phantom{Datum }Ort, Datum\hfill Unterschrift\hspace{15mm}

\newpage

%%%----------------------------------------------------------
%%% Preface
%%%----------------------------------------------------------
\section*{Kurzfassung}\thispagestyle{empty}
Das HTML5 Canvas-Element und die Skriptsprache JavaScript ermöglichen es, den
Browser für Multimedia-Anwendungen und Spiele zu nutzen und die
Internet-Entwicklung unabhängig von Flash zu machen. In Zukunft könnte das
bedeuten, dass Flash-Anwendungen obsolet werden könnten. Im Zuge der Master
Thesis sollen Faktoren wie unter anderem die Möglichkeiten, Performance,
Cross-Plattformkompatibilität und mobile Verwendbarkeit von HTML5 mit Flash
verglichen werden. Die Analysen sollen anhand praktischer Umsetzungen im Bereich
Multimedia- und Spiele-Entwicklung durchgeführt werden. Als Ergebnis soll
ersichtlich sein ob das Canvas-Element und JavaScript die Flash-Entwicklung
ersetzen kann.
\\ \vfill
% Bitte 3-5 deutsche Schlagwörter eingeben, die die Arbeit charakterisieren:
\paragraph*{Schlagwörter:} Webapplikation, Web-Entwicklung, Web Standards,
HTML5, Adobe Flash
\section*{Abstract}\thispagestyle{empty}
The new features of HTML5 offer new possibilities for developers. The new
canvas element in combination with JavaScript can be used to develop
multimedia applications and games for the browser. The result is the
progress of the internet independent from Adobe Flash, which possibly means
that Adobe Flash can be replaced by HTML5.
\newline\newline
This paper compares the core technologies of HTML5 and
Adobe Flash, pointing out how HTML5 and Adobe Flash can
be used at present and which of the two technologies has a greater potential
for the future. On the basis of practical implementations, the advantages and
disadvantages of both competitors are shown and rough guidelines are be
given for the choice of which technology can be preferably used for which
project.
\newline\newline
Analyses based on practical implementations in the field of multimedia and
games include a discussion of the results and future perspectives of the two
technologies.
\\ \vfill
% Bitte 3-5 englische Keywords eingeben, die die Arbeit charakterisieren:
\paragraph*{Keywords:} Web-Entwicklung, Web Standards, HTML5,
Adobe Flash, CSS, JavaScript, ActionScript
\section*{Danksagung}\thispagestyle{empty}
--------DANKSAGUNG-------

\tableofcontents\thispagestyle{empty}
\newpage

\setcounter{page}{1}

%%%----------------------------------------------------------
%%% Hauptteil
%%%----------------------------------------------------------

% =Einleitung
\chapter{Einleitung}
% -Problemstellung
\section{Problemstellung}
Die Entwickler von Flash (früher Macromedia, jetzt Adobe) haben schon früh das
Potential von ihrem Flash Player erkannt, um plattformübergreifend Video- und
Audioinhalte im Web bereitzustellen. Ein großer Vorteil war, dass Inhalte für
Flash lediglich einmal kodiert werden mussten. Tausende von Seiten die sich für
Flash als Streamingplattform von Multimediainhalten entschieden haben,
bestätigten, bis heute, das Potential.
\newline\newline
Mit dem Erscheinen von Apples iPhone und iPod touch im Jahr 2007 und der
folgenden Entscheidung Flash auf diesen Geräten nicht zu unterstützen, mussten
Webseitenbetreiber darauf reagieren. Viele boten Video/Audio-Streams an, die
direkt im mobilen Safari Browser wiedergegeben werden konnten. Durch die
Verwendung der H.264 Kodierung konnten die Inhalte auch über den Flash-Player
(sofern vom verwendeten Gerät unterstützt) abgespielt werden. Dadurch mussten
Inhalte weiterhin nur einmal kodiert werden um mit möglichst vielen Plattformen
kompatibel zu sein.
\newline\newline
Die Entwickler der HTML5 Spezifikation (und u.a. Apple) sind der Meinung, dass
Browser Audio und Video nativ unterstützen sollten ohne dabei auf etwaige
Plugins zurückgreifen zu müssen. Dies hat den Vorteil, dass es keine
herstellerspezifischen Einschränkungen (wie es bei Flash der Fall ist) für den
Entwickler gibt, was er mit den Inhalten anstellt, nachdem sie in die Website
eingebettet wurden. Mittels CSS und JavaScript ist es jederzeit direkt möglich
dargestellte HTML Elemente zu manipulieren.
\newline\newline
Um ein Bild für eine Webanwendung zu erstellen wird üblicherweise eine
Grafiksoftware genutzt und anschließend in die Anwendung eingebettet. Für
Animationen wird vorwiegend Flash verwendet. Mit dem in HTML5 verfügbaren Canvas
Element können Entwickler Bilder und Animationen direkt im Browser mittels
JavaScript generieren. Mit Canvas ist es möglich einfache bis komplexe Formen,
Graphen und auch Diagramme zu erstellen, ohne auf diverse Bibliotheken, Flash
oder ein anderes Plugin zurückgreifen zu müssen.
\newline\newline
Mit der Veröffentlichung und der Implementierung von HTML5 Features in die
neuesten Brower eröffneten sich neue Möglichkeiten um Animationen, Spiele und
Entertainment-Applikationen für das Web zu erstellen. Aufgrund vieler
auftretender Fragen ist sich kein Entwickler wirklich klar ob HTML5 den
Platzhirschen Flash verdrängen kann:
\begin{itemize}
	\item Wo liegen die Vorteile und Nachteile beider Technologien?
	\item Ist HTML5 die Zukunft und löst Flash ab?
	\item Welche Technologie ist einfacher zu erlernen?
	\item Was kann HTML5 was Flash nicht kann (und visa-versa)?
	\item Lohnt es sich noch Flash zu erlernen?
	\item Gibt es Entwicklungsumgebungen die das Programmieren erleichtern?
\end{itemize}
% -Motivation
\section{Motivation}
--------MOTIVATION-------
% -Zielsetzung
\section{Zielsetzung}
Das Ziel dieser Arbeit ist es herauszufinden, ob die kommenden Webstandards
HTML5 und CSS3, im speziellen das Canvas Element, in Zukunft das proprietäre
Webformat Flash samt ActionScript ablösen können. Mittels der Ergebnisse sollen
Richtlinien erstellt werden, anhand derer Web-Entwickler leicht ermitteln
können, in welchen Fällen die eine Technologie der anderen vorzuziehen ist und
aus welchen Gründen.

% -Aufbau der Arbeit
\section{Aufbau der Arbeit}
Kapitel 2 bietet einen Überblick und eine Einführung in die verwendeten
Technologien, die für das Verständnis der Inhalte dieser Arbeit notwendig sind.
\newline\newline
Kapitel 3 
\newline\newline
Kapitel 4
% =Begriffserklaerungen
\chapter{Begriffserklärungen}

% Marc Frydetzki - Flash und HTML5 im Vergleich
Bevor in den nächsten Kapitel die zwei Technologien Flash und HTML5
vorgestellt und im Anschluss anhand einer prototypischen Anwendung
verglichen werden, wird in diesem Kapitel zunächst die Geschichte des
Internets vorgestellt. Beide Technologien nehmen einen großen Stellenwert
in den Diskussionen um dieses Medium ein, wodurch diese Informationen
grundlegend erforderlich sind. Im Anschluss werden die Arbeitsgruppen die
sich mit HTML5 beschäftigen und Begriffe der Thematik Web 2.0 erläutert.

\section{Die Geschichte des Internets}

Der Name Internet kommt von dem englischen Begriff ``Interconnected set of
networks'', was soviel bedeutet wie "`miteinander verbundene Netzwerke"'.
Das Internet ist heute ein weltumspannendes Netz von vielen einzelnen
Computernetzwerken.
Durch die Erfindung des elektrischen Stroms im Jahre 1730 war der
Grundstein für die elektrische Telegrafie gelegt, welche Samuel Thomas
von Soemmerring im Jahr 1809 erfand.
Die Erfindung der elektromagnetischen Induktion im Jahr 1832 durch Michael
Faradey führte zu den ersten Versuchen mit einem elektromagnetischen
Telegrafen, wodurch 1833 die erste telegrafische Nachrichtenübertragung
gelang.
Ab diesem Zeitpunkt befassten sich mehrere Personen mit der Erforschung und
Entwischlung von Apparaturen zur Übertragung von akustischen Signalen.
Schlussendlich war es Alexander Graham Bell - im Jahr 1876 -, der die
Fertigkeiten zusammen brachte, das Telefon über eine Versuchsapparatur
hinaus als Gesamtsystem zur Marktreife zu bringen.
Die Datenübertragung mittels Impulsen, die durch Kabel übertragen werden,
stellt bis heute das wichtigste Transportmittel von Informationen dar.
Aufgeschreckt durch \emph{Sputnik}, dem ersten von der UdSSR ins All
beförderten geostationären Satellit, gründete das Verteidigungsministerium
der USA die``Advanced Research Projects Agency''. Die ARPA hatte die Aufgabe
neue Technologien im Bereich der Datenübertragung und Kommunikation zu
entwickeln.
Das führte Ende der 1950er Jahre, durch J. C. R. Licklider und seinem
Forschungsteam, zur Entwicklung des ersten ``Time-Sharing-Systems'' der Welt.
Dabei war ein Zentralrechner mit sternförmig angeschlossenen Terminals
verbunden, wodurch mehrere Nutzer gleichzeitig dessen Rechenleistung nutzen
konnten. Damit wurde es ermöglicht große geografische Distanzen mit Terminals
eines Herstellers zu überwinden, allerdings war die Anzahl der Anschlüsse
begrenzt. Für militärische Zwecke eignete sich dieser Netzaufbau jedoch nicht,
weil eine Störung des Zentralrechner den Ausfall des gesamten Netzes bedeutete.
Ein weiterführender Ansatz war das dezentrale Netzwerk, mit dem sich Paul
Baran beschäftigte. Hierbei wurden mehrere Zentralrechner, die jeweils über
die sternförmig angeschlossenen Terminals verfügten, miteinander verbunden.
Fiel ein Zentralrechner aus, war das Netzwerk im Ganzen war geschwächt, aber
immer noch in Teilen nutzbar. Um bei einem Netzwerkschaden so wenige Terminals
wie möglich zu verlieren, entwickelte Pauk Barann das ``Distributed Network''
bei dem der Zentralrechner überflüssig wurde. Jedes Terminal, das mit dem
Netzwerk verbunden war, hat alle Funktionen an Board, die für die
Kommunikation im Netzwerk notwendig sind und galt somit als Computer. Ein
weiterer Vorteil mehrerer Verbindungen zu einem möglichen Zielcomputer
führte zur Entwicklung des ``Packet-Switching'', wobei eine Datei nicht mehr
als Ganzes übertragen, sondern in viele Datenpakete gleicher Größe zerteilt
und einzeln transportiert wird. Staus werden somit vermieden, denn pro Paket
kann die Route neu berechnet werden und muss bei einer fehlerhaften
Übertragung nicht noch einmal als Ganzes übertragen werden. Diese intelligente
Übertragung beschleunigte und entlastete das Netzwerk wesentlich. 1965 wurde
diese Art der Netzkommunikation zum ersten Mal eingesetzt. 1966 begann
die Planung zur Vernetzung aller über das Land verteilten Computerzentren
der ARPA nach Paul Barans ``Distributed Network''. Gegen Ende 1969 wurden
nacheinander die großen Computerzentren verbunden.
Die angestrebte Unabhängigkeit von Herstellern und Betriebssystem führte zur
Entwicklung des ``Interface Message Processor'', kurz IMP, durch Wesley Clark.
Der IMP ist ein Minicomputer mit einem einheitlichen Netzwerkprogramm, der an
jeden Großrechner angeschlossen wurde unf dür die Datenübermittlung zuständig
war. Durch diese Verbindung der großen Computerzentren der ARPA enstand das
ARPANET und zum ersten Mal waren Rechner verschiedenster Art miteinander
verbunden. Die ersten Anwendungen im ARPANET waren ``Telnet'' und ``FTP''
bevor Ray Tomlinson 1972 eine Software zum Verschicken und Empfangen
elektronischer Post veröffentlichte. Der E-Mail-Dienst war geboren und
führte zu einem sprunghaften Anstieg der Nutzerzahlen. Weitere Verbreitung
erfuhr das ARPANET durch Vorführung auf der International Conference on
Computer Communications und Weitergabe des Wissens unter anderem an die NASA,
die Air Force und Universitäten. Darauf folgte das ALOHANET, das
Forschungsstationen auf Hawaii vernetzte und sich aufgrund störungsanfälliger
Telefonleitungen zum PRNET (Packet Radio Network) weiterentwickelte, in dem
Daten per Radiowellen übertragen wurden. Als weitere Kommunikationsmöglichkeit
wurde 1973 das SATNET (Datelliten Network) entwickelt. Die Nachteile der
verschiedenen Datenübertragungsarten führten schließlich 1974 Vinton Cerf und
Bob Kahn zur Entwicklung eines umfassenden Netzwerkprotokolls, dass sich auf
ein einheitliches Datenformat und eine einheitliche Verbindungsmethode
beschränkte, dem Transmission Control Protocol / Internet Protocol, kurz
TCP/IP. 1983 wurde das TCP/IP zum Standard erklärt und kommt noch heute in der
Form zum Einsatz. Im selben Jahr entstand das MILNET (Military Network) in
das der komplette militärische Bereich des ARPANET verlagert wurde. Das
ARPANET diente ab dem Zeitpunkt nur noch rein zivilen Gruppen. Dieser
Schritt war notwendig geworden, weil sich immer mehr internationale
Netzwerke anschlossen und somit die Unabhängigkeit des Militärs in Gefahr
geriet. Bis 1985 wurden die bestehenden Netzwerke vorrangig von kleineren und
größeren Nutzergruppen zur Kommunikation und zum Austausch ihrer akademischen
Forschungsthemen genutzt.
1985 gab es einen entscheidende Änderung. Die Netzwerkzugänge des 1984
gegründeten JANET (Joint Academic Network) und des NSFNET (National Science
Foundation Network) wurden für alle Benutzer -- egal aus welchem
Forschungsbereich -- geöffnet. Dadurch erlangte das NSFNET immer mehr an
Beliebtheit und übernahm schließlich 1990 die Funktion des ARPANET. In
diesem Schritt wurde das ARPANET eingestellt. Im selben Jahr schlossen sich
die Netzwerke von Canada, Dänemark, Finnland, Frankreich, Island, Norwegen und
Schweden an das NSFNET an. 1991 folgten dann Deutschland, Japan, Niederlande
und das Vereinigte Königreich.
Der Dienst der in der heutigen Zeit von vielen Menschen mit dem Internet
gleichgesetzt wird ist das World Wide Web. Es ist 1990 aus einer
Weiterentwicklung eines etwas älteren Projekts von Tim Berners-Lee -- er war
in den 80er und 90er Jahren Informatiker am Hochenergieforschungszentrum
CERN -- hervorgegangen und führte zu einer rasant ansteigenden Beliebtheit
des Internet.

\section{Browser}
Ein Browser ist ein spezielles Programm, das zur Darstellung von Inhalten aus
dem Internet (World Wide Web) genutzt werden kann. Er übersetzt die im
Quellcode verwendeten Befehle in das entsprechende Aussehen. Browser sind
die Benutzeroberfläche für Webanwendungen, aber auch andere Daten und
Dokumente können betrachtet werden. Es gibt mehrere Hersteller, die die
Konventionen des W3C unterschiedlich umsetzen. Diese auseinander gehenden
Herangehensweisen führen dazu, dass Webseiten in verschiedenen Browsern
unterschiedlich dargestellt werden. In unterschiedlichen Versionen eines
Browsers können Konventionen unterschiedlich umgesetzt worden sein, was die
versionsübergreifende Entwicklung für Programmierer zu einer enormen
Herausforderung macht. Zu den Merkmalen eines guten Browsers zählen
Schnelligkeit, Leistungsbedarf auf dem Rechner des Anwenders und die
Umsetzung der neuesten Innovationen im World Wide Web.

\section{W3C}
Das W3C ist ein internationales Konsortiom, in dem Mitgliedsorganisationen, ein
fest angestelltes Team und die Öffentlichkeit gemeinsam daran arbeiten,
Web-Standards und Richtlinien zu entwickeln -- daher der Name World Wide Web
Consortium, kurz W3C. Gegründet wurde es am 1. Oktober 1994 von Tim
Berners-Lee, dem Erfinder des World Wide Web. Die Entstehung des W3C ist eng
mit der Entstehung des World Wide Web verbunden. Die große Gefahr von
Inkonsistenten in der Nutzung vorliegender Technologien könnte zu unwirksamen
Verknüpfungen führen und das galt es möglichst zu verhindern. Die
Grundherangehensweise an neue Technologien oder Teilaspekte besteht darin, den
kleinsten gemeinsamen Nenner zu finden und diesen zu einer Spezifikation zu
verarbeiten, so dass diese von allen Mitgliedsorganisationen unterstützt wird.

\section{WHATWG}
Die Web Hypertext Application Working Group, kurz WHATWG, ist eine
Arbeitsgruppe, die von mehreren Unternehmen, darunter unter anderem Mozilla
Foundation, Opera Software ASA und Apple Inc., betrieben wird und deren Ziel
die Entwicklung neuer Technologien zur einfacheren Erstellung von
Internetanwendungen ist. Die WHATWG wurde 2004 gegründet, weil aus Sicht der
Browser-Hersteller das W3C die Entwicklung von HTML -- zu Gunsten von XHTML
-- vernachlässigte und auf die Bedürfnisse von Programmierern zu wenig
einging. Als Unterschied zu den Bestrebungen der W3C muss festgehalten werden,
dass die WHATWG keinen festen Standard als Zielsetzung verfolgt. Vielmehr soll
HTML weiterentwickelt werden, ohne aber auf einen definierten Stand hin zu
arbeiten. Anders das W3C, dessen aktuelle Bestrebungen ganz klar auf HTML5
ausgelegt sind.

\section{Plattformunabhängigkeit}
Plattformunabhängigkeit bedeutet bei Software, dass sie ohne extra Aufwand und
ohne Hilfssoftware auf unterschiedlichen Systemen ausgeführt werden kann.
Genauer betrachtet ist darunter die Eigenschaft eines Programms zu verstehen,
auf unterschiedlichen Computersystemen mit Unterschieden in Architektur,
Prozessor, Compiler, Betriebssystem und weiteren Dienstprogrammen, die zur
Übersetzung oder Ausführung notwendig sind, lauffähig zu sein.

\section{Rich Internet Application}
Der Begriff ``Rich Internet Application'' kurz RIA (reichhaltige
Internetanwendung) wurde von Macromedia im Jahre 2002 erstmalig verwendet.
Damalige Internetangebote entwickelten sich zunehmend vom reinen statischen
Wiedergeben von Inhalten zu vom Nutzer manipulierbaren dynamischen Inhalten,
wodurch der Begriff \emph{Web 2.0} entstand. Der Begriff RIA
bezeichnet ein Konzept und basiert meist auf mehreren Web-Technologien
wie zum Beispiel HTML, CSS, JavaScript und AJAX oder Flash und PHP. Um als RIA
zu gelten, muss eine Anwendung über das Internet verfügbar sein. Allerdings
muss sie nicht zwangsläufig im Browser ausgeführt werden, sondern kann
auch als Desktopanwendung zum Einsatz kommen. RIAs im heutigen Sinn bieten dem
Anwender Möglichkeiten, die früher Desktopanwendungen zur Verfügungen stellten.
Die Individualisierung und Vereinfachung der Internetnutzung spielt eine
entscheidende Rolle bei der Entwicklung neuer Anwendungen. Unter Interaktion
versteht man beispielsweise das Ändern der Gestaltung, das Erzeugen und
Verwalten von Daten, Drag\&Drop-Funktionalitäten, sowie die Verwendung von
Tastenkürzeln.
%% =Überblick
\chapter{Überblick}
% -HTML5 Geschichte
\section{Die Geschichte von HTML5}
HTML ist die allgegenwärtige Auszeichnungssprache des World Wide Webs. Durch den geschickten Einsatz der paar Elemente (fortan: Tags) die die Sprache unterstützt, konnten erstaunlich viele unterschiedliche Netzwerke von verlinkten Dokumente erstellt werden. Von bekannten Seiten wie Amazon, Ebay und Wikipedia bis hin zu personalisierten Blogs und Webseiten die sich auf Katzenbilder spezialisiert haben.
HTML5 ist die aktuellste Version dieser Auszeichnungssprache. Obwohl diese
Version die bisher umfangreichsten Änderungen mit sich bringt, ist es nicht die
erste Aktualisierung von HTML.
Sir Tim Berners-Lee zeichnet sich für die Entwicklung von HTML und damit dem
Beginn des Internet verantwortlich. 1991 veröffentlichte er ein Dokument mit dem
Titel ``HTML Tags'', in dem er weniger als zwei Dutzend Elemente, die für das
Schreiben von Webseiten genutzt werden konnten, vorschlug.
Die Verwendung von Tags (Wörter umgeben von eckigen Klammern: z.B. <html>) war
nicht Berners-Lees eigene Errungenschaft. Tags wurden bereits in der SGML
(Standard Generalized Markup Language) verwendet.
Anstatt einen komplett neuen Standard zu erfinden, erkannte Berners-Lee die
Vorteile, bereits existierende Standards weiterzuentwickeln - ein Trend der auch
in der Entwicklung der neuen HTML5 Spezifikation erkennbar ist.

\subsection{Von der IETF zur W3C: Der Weg zu HTML4}
Die erste offizielle Spezifikation war HTML 2.0, veröffentlicht durch die IETF,
die ``Internet Engineering Task Force''. Viele der neuen Features dieser
Spezifikation basierten dabei auf bereits veröffentlichte Implementierungen. So
bot der im Jahr 1994 marktführende Web Browser Mosaic Autoren von Webseiten
bereits die Möglichkeit Bilder in Dokumente mittels eines <img>-Tags
einzubinden. Das <img>-Tag wurde in die HTML 2.0 Spezifikation inkludiert.
Die IETF wurde allmählich durch die W3C, dem World Wide Web Consortium, ersetzt.
Folgende Aktualisierungen des HTML Standards wurden auf \url{http://www.w3.org}
veröffentlicht. In der zweiten Hälfte der Neunziger wurde der Standard mehrmals
überarbeitet bis 1999 HTML 4.01 veröffentlicht wurde.

\subsection{XHTML 1: Die Vermischung von HTML und XML}
Das nächste Update von HTML 4.01 trug den Namen XHTML 1.0. Das X steht dabei für
``eXtensible''. Die Spezifikation von XHTML 1.0 war ident zu der von HTML 4.01,
somit wurden keine neuen Elemente und Attribute eingeführt. Die Spezifikationen
unterschieden sich nur durch die zu verwendende Syntax zum Schreiben von
Dokumenten. Während Autoren von HTML Dokumenten kaum Einschränkungen im
Schreibstil von Elementen und Attributen hatten, setzte XHTML es voraus, dass
die Regeln von XML, einer strikteren Auszeichnungssprache auf der viele
Technologien des W3C basierten, eingehalten werden.
Aufgrund der strikteren Regeln einigten sich Autoren auf einen gängigen
Schreibstil, der auch unter dem HTML 4.01 Standard Verwendung fand. Während Tags
früher in Kleinbuchstaben, Großbuchstaben oder einer Mischung aus beidem
geschrieben werden konnten, verlangte ein valides XHTML Dokument, dass alle
Elemente und Attribute klein geschrieben werden.
Während XHTML 1.0 noch auf HTML basierte und lediglich die strikteren Regeln von
XML verwendete, war die XHTML 1.1 Spezifikation reines XML. Dieser Umstand
führte zu schwerwiegenden Problemen. XHTML 1.1 Dokumente konnten nicht mehr
unter dem Mime-Type \textit{text/html} definiert werden. Der bis dato
bekannteste Webbrowser Internet Explorer konnte jedoch Dokumente die mit einem
XML Mime-Type publiziert wurden, nicht darstellen.

\subsection{XHTML 2}
Das W3C war mit der vierten Version von HTML der Meinung, dass der HTML basierte
Ansatz seinen Zenith erreicht hat und setzten für die zukünftigen Version
vollständig auf XML. Trotz der fast identen Namen von XHTML 1 und XHTML 2,
konnten die Unterschiede zwischen den beiden Spezifikationen nicht größer sein.
Anders als XHTML 1, war XHTML 2 nicht abwärtskompatibel mit bereits
existierenden Webinhalten oder vorangegangenen HTML Versionen. Es sollte ein
vollkommen neuer Standard werden und es zeigte sich, dass dieser Weg nicht mit
Erfolg gekrönt sein würde.

\subsection{Die Spaltung: WHATWG}
Innerhalb des W3C bildete sich eine Gruppierung die gegen die Ansätze
rebellierten. Namenhafte Vertreter von Opera, Apple und Mozilla waren mit der
eingeschlagenen Entwicklungrichtung des W3Cs nicht zufrieden. Ihnen war es
wichtiger mehr Aufmerksamkeit auf Formate zur Entwicklung von Webapplikationen
zu richten.
Bei einem Workshop in 2004 schlug Ian Hickson, der zu der Zeit bei Opera
arbeitete, eine Weiterentwicklung von HTML vor, mit der es möglich sein soll
Anwendungen zu entwickeln. Der Vorschlag wurde durch die W3C abgelehnt.
Unzufrieden mit der Entscheidung bildeten die Rebellen eine eigene Gruppe: Die
Web Hypertext Application Technology Working Group, oder kurz WHATWG.

\subsection{Der Weg von Web Apps 1.0 zu HTML5}
Von Beginn an operierte die WHATWG anders als das W3C. Während das W3C einen
kosensorientierten Ansatz anwendet, Themen werden vorgetragen, es wird
diskutiert und abgestimmt, wird auch bei der WHATWG diskutiert und abgestimmt,
allerdings liegt die finale Entscheidung, was in die Spezifikation komm und was
nicht, beim Editor. Der Editor ist Ian Hickson.
Der W3C Ansatz klingt demokratisch und fair, allerdings führt dieser auch dazu,
dass der Prozess stark verlangsamt wird. Bei der WHATWG hat jeder die
Möglichkeit mitzuwirken, da die letzte Entscheidung aber beim Editor liegt,
entwickelt sich alles schneller.
Schon zu Beginn der Gründung von WHATWG wurden die Projekte in zwei große
Spezifikationen aufgeteilt: Web Forms 2.0 und Web Apps 1.0. Beide
Spezifikationen sollen die bisherige HTML Spezifikation erweitern. Mit
fortschreitender Entwicklung wurden alle Spezifikationen in eine einzige
implementiert und in HTML5 umbenannt.

\subsection{Die Wiedervereinigung}
Während HTML5 von der WHATWG weiter entwickelt wurde, setzte das W3C die
Entwicklung der XHTML 2 Spezifikation fort. Erst im Oktober 2006, schrieb Sir
Tim Berners-Lee einen Blog-Post in dem er zugab, dass der Versuch, das Internet
von HTML auf XML zu übertragen, nicht funktionieren wird.
Nur wenige Monate später entschied sich die W3C eine neue HTML Working Group zu
bilden. Anstatt komplett von Anfang an zu beginnen, nutzen sie glücklicherweise
die bisherige Arbeit der WHATWG als Basis für zukünftige HTML Spezifikationen.
Dieses hin- und her führte zu einer schwer überschaubaren Situation. Das W3C
arbeitete parallel an den beiden unterschiedlichen, nicht kompatiblen
Spezifikationen: XHTML 2 und HTML 5 (mit Leerzeichen). Währenddessen arbeitete
die WHATWG, an der Spezifikation für HTML5 (ohne Leerzeichen) auf der die Arbeit
der W3C aufbaut.

\subsection{HTML5: 2012 und 2022}
Heute ist der aktuelle Stand der HTML5 Spezifikation nicht mehr so
undurchsichtig wie er früher war, allerdings gibt es noch immer offene Fragen.

Es gibt noch immer zwei Gruppen die an HTML5 arbeiten.
Die wohl wichtigste Frage für Webentwickler ist "`Wann können wir es nutzen?"'.
In einem Interview gab Ian Hickson an, dass HTML5 frühestens 2022 den ``proposed
recommendation'' Status erreichen wird. Klingt nach einer langen Wartezeit,
allerdings bedeutet ``proposed recommendation'', dass die HTML5 Spezifikation
zwei mal komplett implementiert werden muss. Als Vergleich: HTML 4 existiert nun
seit über einem Jahrzehnt und hat noch nicht die gesetzten Features erreicht.
Wenn man den Umfang der Spezifikation betrachtet, klingt dieses Datum hoch
gesteckt. Browserhersteller sind nach wie vor nicht dafür bekannt, dass
existierende Standards so schnell wie möglich implementiert werden. Der Internet
Explorer benötigte mehr als ein Jahrzehnt um das \textit{abbr}-Element richtig
darstellen zu können.

Für Webentwickler war das Jahr 2012 wesentlich wichtiger. 2012 erreichte die
HTML5 Spezifikation den ``candidate recommendation'' Status, der gleichbedeutend
ist mit "`fertig und abgeschlossen"'.
Allerdings reicht das alleine leider auch nicht aus. Wirklich entscheident ist
es wann Webbrowser den neuen Standard unterstützen. Schon die Veröffentlichung
des CSS 2.1 Standards zeigte, dass man nicht auf die Fertigstellung der
Spezifikationen warten sollte, sondern, wenn möglich, die Features nutzen sollte
sobald es möglich ist. Das selbe gilt auch für HTML5. Sobald Webbrowser
bestimmte Features der Spezifikation unterstützen können diese auch jederzeit
verwendet werden.
Man darf nicht vergessen, dass HTML5 keine komplett neu entwickelte Sprache ist.
Im Sinne der HTML Spezifikation ist es eher eine Evolution als eine Revolution.
Da HTML5 auf den früheren Versionen aufsetzt und irgendeine Version des HTML
Standards zur Erstellung von Webseiten genutzt wird, wird bereits HTML5 genutzt.


% \section{Geschichte}
% \subsection{Der Begründer des World Wide Web}
% Die Geschichte des Internets geht zurück bis in die frühen neuziger als Tim
%Berners-Lee das World Wide Web entwickelte mit der Hyper Text Markup Language
%(HTML) als Auszeichnungssprache. Während die Elemente (auch Tags genannt) zur
%Auszeichnung von Inhalten in strukturierte Teile (Paragraphen, Überschriften,
%Listen, etc.) dabei auf der Standard Generalized Markup Language (SGML), eine
%international anerkannte Methode um Texte zu strukturieren, basierte, war die
%Idee des Hypertext Links von Tim Berners-Lee.
%Ein Jahr später veröffentlichte Berners-Lee die erste Version seines Browsers.
%Während dieser Zeit waren die Möglichkeiten von HTML stark eingeschränkt und
%Entwickler konnten lediglich einfache Textinhalte im Web veröffentlichen. Mit
%dem Erscheinen von HTML+, welche von Dave Ragget von Hewlett-Packards Labs
%entwickelt wurde, gab es die Möglichkeit mittels IMG Element auch Bilder in
%Webseiten zu inkludieren

%Versuche um weitere Verbesserungen zur HTML Spezifikation hinzuzufügen, führte
%zur Veröffentlichung einer neuen Version, HTML 2.0. Sie basierte vollständig auf
%der alten Version wurde jedoch mit weiteren Funktionen und der Document Type
%Definition (DTD) erweitert. Die Funktionen von HTML 2.0 wurden zum Standard für
%alle Webbrowser. Der Wunsch nach Personalisierung, gestalterischer Freiheit und
%einem zeitgemäßen Look and Feel führte zu HTML 3.0 und der Einführung der
%Cascading Style Sheets.
%
%Probleme entstanden erst als Netscape sich entschied eigene Tags und Attribute
%in ihre Browser zu implementieren. Entwickler mussten Tags speziell für Netscape
%anpassen. Netscape sah jedoch ein, dass ihr Versuch nicht zielführend ist und
%gab ihe Erweiterungen auf. Mit dem steigenden Interesse am Internet
%veröffentlichte Microsoft die erste Version des Internet Explorers.
%
%Mit dem schnellen Fortschreiten der Entwicklung des HTML Standards wurde HTML
%4.0 veröffentlicht. Diese Version enthielt weitere Verbesserungen und
%Möglichkeiten in der Entwicklung von Webseiten. Es gab mehr Möglichkeiten um
%Multimedia in Webseiten einzusetzen, nutzen von Skiptsprachen, verbesserte
%Druckeigenschaften, Style Sheets und Optimierungen für Menschen mit
%Behinderungen.

%IMAGE Timeline of Web Technologies
% http://www.instantshift.com/2012/07/20/the-evolution-of-html5-infographic/
% publishing multimedia on the web.pdf

\section{Rich Media}
\subsection{Audio}
MP3 stellt das allgegenwärtige Format zur Kodierung von Audio-Dateien dar. Um sich diese
Audio-Dateien auch anhören zu können waren propritäre Technologien notwendig. Damit
wurde der Flash Player allgegenwärtig.
\newline\newline
Mit HTML5 bietet sich eine neue Technologie an die versucht den Platzhirschen Flash zu
verdrängen.
Das Einbinden einer Audio-Datei in ein HTML5-Dokument ist äußerst simpel:
\begin{verbatim}
<audio src="sample.mp3">
</audio>
\end{verbatim}
Das Audio Element bietet mehrer Attribute zum Steuern der Audioausgabe an.
Das Attribut ``autoplay'' sorgt dafür, dass die Audio-Datei nach vollständigem Laden
sofort gestartet wird. ``loop'' hingegen, kümmert sich darum, dass das die Audio-Datei
in einer Endlosschleife abgespielt wird.
\begin{verbatim}
<audio src="sample.mp3" autoplay loop>
</audio>
\end{verbatim}
Eine Besonderheit dieser Attribute ist, dass sie im Gegensatz zu den herkömmlichen HTML-Attributen,
keinen Wert zugewiesen bekommen. Das liegt daran, dass es sich bei diesen Attributen bereits um
``Boolean''-Attribute handelt. Das zuweisen eines Wertes, z.B. autoplay="no", wird an der Ausführung
des Attributes nichts ändern - entweder die Attribute werden angegeben und ausgeführt oder
eben nicht.

\subsubsection{Audio API}
Das Boolean-Attribut controls sorgt dafür, dass der Browser die nativen Kontrollmöglichkeiten
anzeigt und den Usern zu Verfügung stellt. Somit ist es sehr schnell möglich die Audiowiedergabe
zu Starten/Stoppen, die Position und die Lautstärke zu verändern.
%Image Audioplayer mit Controls
Über JavaScript ist es möglich die Audio API zuzugreifen und somit Methoden für Play, Pause und
Eigenschaften wie Lautstärke zu nutzen. Ein einfaches Beispiel mit Button Elementen und
Inline Event Handler:
\begin{verbatim}
<audio id="player" src="sample.mp3">
</audio>
<div>
<button onclick="document.getElementById('player').play()">
Play
</button>
<button onclick="document.getElementById('player').pause()">
Pause
</button>
<button onclick="document.getElementById('player').volume += 0.1">
Volume Up
</button>
<button onclick="document.getElementById('player').volume -= 0.1">
Volume Down
</button>
</div>
\end{verbatim}

\subsubsection{Audioformate}
Obwohl das Audio Element einen sehr guten Eindruck macht gibt es einen Wehrmutstropfen,
und dieser liegt nicht in der Spezifikation des Elements. Das Problem ist die breite Fragmentierung
von Audiocodecs. Leider gibt es bezüglich der zu verwendeten Codecs unterschiedliche Meinungen
unter den Browserherstellern. Während heutzutage das MP3 Format allgegenwärtig ist,
ist es nach wie vor kein offenes Format. Die Folge ist, dass MP3 Dateien von Applikationen nur
dann dekodiert werden können, wenn für die entsprechenden Patentrechte bezahlt wird.
Für Giganten wie Apple oder Adobe stellt das kein gröberes Problem dar, allerdings erschwert es die
Arbeit von kleineren Unternehmen und Open-Source Organisationen. Infolgedessen gibt
Apples Safari ohne Probleme MP3 Datei wieder während Mozillas Firefox daran scheitert.
\newline\newline
Natürlich existieren weitere Audioformate. Der Vobis Codec - auch bekannt unter der Datei-Endung
.ogg - ist zum Beispiel nicht patentiert. Firefox unterstützt den Codec aber Safari nicht.
\newline\newline
Glücklicherweise ist es nicht notwendig eine fixe Entscheidung bei der Auswahl des Codecs zu treffen.
Anstatt das src Attribut im sich öffnenden <audio> Tag zu nutzen, können mehrere Dateiformate
über source Elemente festgelegt werden:
\begin{verbatim}
<audio controls>
<source src="sample.ogg">
<source src="sample.mp3">
</audio>
\end{verbatim}
Ein Browser der Ogg Vorbis Datei wiedergeben kann wird sich für die weiteren source Elemente
nicht mehr interessieren. Ein Browser der MP3 wiedergeben kann aber Ogg Vorbis nicht wird
einfach das erste Source Element überspringen und die Datei im zweiten wiedergeben.
Zusätzliche Hilfe bei der Entscheidung bietet die Deklarierung des entsprechenden
Mime Types der Audio Datei:
\begin{verbatim}
<audio controls>
<source src="sample.ogg" type="audio/ogg">
<source src="sample.mp3" type="audio/mpeg">
</audio>
\end{verbatim}
Um die Möglichkeiten vom Audio Element vollständig ausnutzen zu können wird angeraten
MP3 und Ogg Vorbis als Codecs zu verwenden.

\subsubsection{Fallback Lösungen}
Auch wenn das festlegen von mehreren source Elementen sehr nützlich ist, muss
bedacht werden, dass Browser existieren die das Audio Element nicht unterstützen.
Internet Explorer und Konsorten müssen auf die altmodische Art auf ein Flashumsetzung
zurückgreifen. Glücklicherweise unterstützt das Audie Element die Nutzung von Flash.
Alles was zwischen den sich öffnenden und schließenden <audio> Tags befindet und
kein source Element ist wird nur dem Browser, der das audio Element nicht unterstützt,
angezeigt:
\begin{verbatim}
<audio controls>
<source src="sample.ogg" type="audio/ogg">
<source src="sample.mp3" type="audio/mpeg">
<object type="application/x-shockwave-flash"
<param name="movie" value="player.swf?soundFile=sample.mp3">
</object>
</audio>
\end{verbatim}
Das object Element bietet zusätzlich eine Möglichkeit an um Fallback Inhalte anzubieten.
So kann zum Beispiel im schlimmsten Fall ein gewöhnlicher Downloadlink angezeigt werden.
\begin{verbatim}
<audio controls>
<source src="sample.ogg" type="audio/ogg">
<source src="sample.mp3" type="audio/mpeg">
<object type="application/x-shockwave-flash"
<param name="movie" value="player.swf?soundFile=sample.mp3">
<a href="sample.mp3">Audio Datei herunterladen</a>
</object>
</audio>
\end{verbatim}
Mit diesem Code werden bereits vier Fallback Ebenen angeboten:
\begin{itemize}
\item{Der Browser unterstützt das Audio Element und den Ogg Vorbis Codec.}
\item{Der Browser unterstützt das Audio Element und den MP3 Codec.}
\item{Der Browser unterstützt das Audio Element nicht, hat aber das Flash Plug-in installiert.}
\item{Der Browser unterstützt das Audio Element nicht und hat kein Flash Plug-in installiert.}
\end{itemize}

\subsection{Video}
Mit dem anstieg der verfügbaren Bandbreite stieg auch das Interesse an Video Inhalten an.
Das Flash-Plugin ist derzeit noch die erste Wahl wenn Videos im Web angeboten werden sollen.
Mit HTML5 könnte sich das ändern.
\newline\newline
Das Video-Element funktioniert genauso wie das Audio-Element. Es unterstützt die
selben optionalen ``autoplay'', ``loop'' und ``preload'' Attribute. Der Speicherort des Videos
kann entweder mittels ``src'' Attribut im Video-Element oder mittels ``source'' Elementen, die
sich verschachtelt zwischen den sich öffnenden und schließenden <video> Tags befinden,
festgelegt werden. Zur Darstellung eines geeigneten User Interfaces kann entweder mittels
``controls'' Attribut dem der Browser die Darstellung übernehmen lassen oder es wird mit
entsprechenden HTML-Elementen, CSS Befehlen und JavaScript ein benutzerdefiniertes erstellt.
\newline\newline
Einer der wesentlichen Unterschiede zwischen dem Audio und Video Element ist, dass Videos
natürlicherweise einen fixen Bereich der Webseite einnehmen werden. Um diesen
Bereich festzulegen müssen im Video Element die entsprechenden Dimensionen definiert werden:
\begin{verbatim}
<video src="sample.mp4" controls width="360" height="240">
</video
\end{verbatim}
Um beim Laden des Videos kann mittels ``poster'' Attribute dem Browser mitgeteilt werden,
währenddessen ein representatives Bild anzuzeigen:
\begin{verbatim}
<video src="sample.mp4" controls width="360" height="240"
poster="sampleimage.jpg">
</video
\end{verbatim}
% Image Placeholder and Dimensions
Der Kampf der rivalisierenden Videoformate stellt sich noch extremer als bei den Audio Formaten dar.
Die am stärksten vertretensten Formate sind das patentierte MP4 und das freie Theora Video. Wie beim
Audio Element muss das Video in mehreren Kodierungen verfügbar und Notfalls eine Fallback Lösung
vorhanden sein.
\begin{verbatim}
<video controls width="360" height="240"
poster="sampleimage.jpg">
<source src="sample.ogv" type="video/ogg">
<source src="sample.mp4" type="video/mp4">
<object type="application/x-shockwave-flash" width="360" height="240"
data="player.swf?file=movie.mp4">
<a href="movie.mp4">Video herunterladen</a>
</object>
</video
\end{verbatim}
Die Autoren der HTML5 Spezifikation hofften auf eine Standardisierung des Videoformates. Allerdings
konnten sich die Browserhersteller, bis heute, nicht auf ein einziges Format einigen.

\subsubsection{Native Unterstützung von Videos}
Die Fähigkeit Videos nativ in Webseiten einzubinden stellt womöglich eine der aufregensten Erweiterung
von HTML dar, seit der Einführung des ``img'' Elements. Giganten wie Google zögern nicht lange und
zeigen bereits ihren Enthusiasmus mit einer auf HTML5 basierenden YouTube-Version:
\url{http://youtube.com/HTML5}
Eines der Probleme von Plug-ins zur Darstellung von Webinhalten ist, dass der Inhalt des Plug-ins
von dem restlichen Inhalt der Webseite geschützt ist (``sandboxed''). Nativen Rich Media Elementen
in HTML haben zur Folge, dass sie ohne Probleme mit den anderen Web-Technologien,
CSS und JavaScript, zusammen arbeitet.
\newline\newline
Das Video Element ist somit nicht nur programmierbar sondern auch stylebar.
% Image Skin & Style
Ein Plug-in bietet derartige Möglichkeiten nicht an.

\subsection{Canvas}

Das Canvas Element ist eine Umgebung zur Erstellung von dynamischen Bildern.
Das Element ist genauso einfach wie das Audio oder Video Element zu verwenden.
Als Attribute werden lediglich Breiten- und Höhenangaben des Canvas angeboten:
\begin{verbatim}
<canvas id="canvas" width="360" height="240">
</canvas>
\end{verbatim}
Alles was sich zwischen den sich öffnenden und schließenden <canvas> Tags
befindet, wird nur Browsern angezeigt, die das canvas Element nicht
unterstützen.
\begin{lstlisting}[language=html]
<canvas id="canvas" width="360" height="240">
<p>Canvas Element wird von ihrem Browser nicht unterstuetzt.</p>
</canvas>
\end{lstlisting}

JavaScript wird verwendet um das Canvas produktiv nutzen zu können.
Um mit dem Canvas arbeite zu können muss immer das entsprechende
Element über ihre ID ausgewählt und der Kontext festgelegt werden.
Kontext bedeutet in diesem Fall welche API genutzt werden soll:
\begin{verbatim}
var canvas = document.getElementById('canvas');
var context = canvas.getContext('2d');
\end{verbatim}
Aktuell stehen nur der 2D und WebGL Kontexts zur Verfügung.
Mit der Auswahl des Kontextes ist das Canvas Element bereit
um für das Zeichnen von Bildern genutzt zu werden.
Die 2D API verfügt über so ziemlich alle Tools die man auch bei einem
Grafik Programm wie Illustrator erwartet: Es können Linien, Konturen, gefüllte
Flächen, Verläufe, Schatten, Formen und Bézier Kurven gezeichnet werden.
Der wesentliche Unterschied besteht allerdings darin, dass kein
grafischen User Interface genutzt wird, sondern alles mit JavaScript
definiert werden muss.

\subsubsection{Mit Code zeichnen}
Um die Farbe einer Kontur bzw. einer Linie zu definieren ist folgender
Code notwendig:
\begin{verbatim}
context.strokeStyle = "\#990000";
\end{verbatim}
Alles was nun auf dem Canvas gezeichnet wird, hat eine rote Kontur.

Die Syntax für die strokeRect Methode sieht dabei folgendermaßen aus:
\begin{verbatim}
strokeRect(left, top, width, height);
\end{verbatim}
Um nun ein rotes Rechteck zeichnen, das 20 Pixel vom linken Rand und
30 Pixel vom oberen Rand des Canvas entfernt ist und 100 Pixel breit und
50 Pixel hoch ist, ist folgender Code notwendig:
\begin{lstlisting}
context.strokeRect(20, 30, 100, 50);
\end{lstlisting}
% Image Rectangle Canvas
Dabei handelt es sich noch um ein sehr einfaches Beispiel. Die 2D API
bietet eine sehr umfangreiche Auswahl an Methoden wie
fillStyle, fillRect, lineWidth, shadowColor und viele mehr.
\newline\newline


\section{Verfügbarkeit}
Seit der offiziellen Präsentation der ersten Entwürfe der erneuerten
Webtechnologie, ist die Begeisterung in der IT Branche groß. Diese Begeisterung
ist besonders an dem Enthusiasmus der Web-Giganten Apple, Google und Mozilla und
der stetigen Implementierung von HTML5 und CSS3 in ihre Browser erkennbar.

\subsection{HTML5, CSS3 und JavaScript: Das Web von morgen}
Graph

\subsection{Verfügbarkeit auf Computer}
Ob und wie weit die Implementierung der neuen Webfeatures in Browsern
fortgeschritten ist, hängt vollständig von den jeweiligen Browserherstellern ab.
Um die Verfügbarkeit auf Computern festzustellen, muss der Marktanteil der
unterschiedlichen Browser mit einbezogen werden.

IMAGE Browser Jahr

Es ist notwendig die kommenden Trends zu bedenken, da die weitere
Implementierung von HTML5 stark von der Popularität des jeweiligen Web-Browsers
abhängt. Unbekanntere Browserhersteßller werden naturgemäß mehr Zeit für die zur
Verfügungstellung von HTML5-Features benötigen als bekannte und beliebte Größen
wie Google oder Mozilla.

IMAGE Browser Statistik Monat zu Monat

Die angeführten Statistiken zeigen auf, dass nicht nur der gegenwärtige
Prozentsatz an Nutzern sondern auch der Trendverlauf darauf schließen lässt,
dass der Internet Explorer tatsächlich nicht mehr zu den marktführenden Browsern
zählt, wodurch auch dessen Bedeutung in der Web-Entwicklung schwindet.
Aufgrund der sich reduzierenden Nutzeranzahl und dem Umstand, dass der Internet
Explorer der Browser mit der am geringsten fortgeschrittensten Implementierung
des HTML5 Spezifikation am Markt ist, ermöglicht vielen Web-Entwicklern, schon
vor der Fertigstellung der Spezifikation, die Nutzung einiger neuen Features.

Microsofts Marktstrategie ihr eigenes Betriebssystem (Windows) ausschließlich
mit dem eigenen Browser auszuliefern ändert kaum etwas an der Statistik. Viele
Benutzer installieren daher zum Beispiel Mozilla Firefox oder Google Chrome
selbstständig.

\subsection{Verfügbarkeit auf mobilen Geräten}

\section{Möglichkeiten}

Folglich werden einige neue Elemente der HTML5 Spezifikation genauer
vorgestellt. Diese wurden auch in der Umsetzung der praktischen Ausarbeitungen
verwendet.


\subsection{Spiele}
\subsection{Animationen}
Mit HTML5 Canvas und den CSS3 Neuerungen gibt es für Entwickler neue APIs um
Grafiken im Web zu zeichnen und diese auch zu animieren. Die APIs bieten
einfache anzuwendente Funktionen an um vordefinierte Formen zu zeichnen, Bilder
zu importieren, die Darstellung von bereits gezeichneten Bildern zu verändern
oder diese auch zu animieren. Im Gegensatz dazu bietet Flash eine komplette IDE
um schnell komplexe Formen zu erstellen, diese anzuzeigen, deren Verhalten zu
steuern und zu animieren.

Ein Charakter lässt sich mit einer Zeichenapplikation um ein vielfaches
einfacher gestalten als mit mehreren Zeilen Code. Flash bietet eine
vollständiges Tool zur Erstellung und Animierung komplexer Grafiken und ist
damit wesentlich effizienter zu nutzen, solange es kein equivalentes Tool für
HTML5 und CSS3 gibt.

Dieses Kapitel soll aufzeigen, ob HTML5 und CSS3 effizient für die Erstellung
von Animationen verwendet werden kann. Anhand bestehender Animationen ...
Anschließend werden die Fragen behandelt wie eine Animation mittels Code
erstellt wird und ob sich der Zeitaufwand für Entwickler auszahlt. Abschließend
wird analysiert ob Webbrowser die Animationen flüssig wiedergeben können.

\subsubsection{CSS3 Spiderman}
Ein kurzer Animationsfilm der mittels HTML5, CSS3 und jQuery erstellt wurde. Die
Animationen sind komplex. So bewegt sich der Hintergrund, es gibt mehrere
Szenen, unterschiedliche Betrachtungswinkel, Bewegungen und Gesichtsausdrücke.
Der Film beinhaltet auch Musik mittels dem HTML5 Audio Element.

Bisher gibt es viele Experimente die mittels CSS3 umgesetzt wurden. Im Vergleich
gibt es nur wenige die auf das Canvas Element und JavaScript setzen.

CSS3 bietet einfache Transformationen (Rotation, Translation, Skalierung) und
reichen für simple Animationen. Flash bietet zusätzlich Formtransformationen,
als Beispiel die allmähliche Veränderung eines Kreises in ein Quadrat. Bisher
ist das mit CSS3 nicht möglich.

CSS3 zeigt im Vergleich zu HTML5 ein größeres Potential um Animationen zu
erstellen. Allerdings ist es schwierig mittels Code komplexere Animationen zu
erstellen. Solange es kein Tool für HTML5 und CSS3 Animationen existiert wird
Flash weiterhin die führende Software in diesem Bereich bleiben.

Performancetechnisch wurde diese Animation ohne jegliche Probleme flüssig in
jedem modernen Browser angezeigt.

\subsection{Entertainment}

\subsubsection{Audio \& Video}
Flash ist bis heut noch das bekannteste propritäte Plugin um Audio- und
Videoinhalte im Internet zur Verfügung zu stellen. Jedoch bietet nun HTML5
eigene Features die diesem Bereich gewidmet und sind ernst zu nehmende
Alternativen zu Flash. Es werden die Audio und Video Spezifikation von HTML5
genauer betrachtet.

\subsubsection{Audio}
Das Audio Element steht bisher in allen aktuellen Webbrowsern zur Verfügung und
ist ähnlich einfach zu nutzen wie das Image Element. Folgender Code zeigt wie
ein Audiofile in eine Webseite eingebunden werden kann:
\begin{verbatim}
<audio src="audio.ogg" controls>
<p>Your browser does not support the audio element.</p>
</audio>
\end{verbatim}
Leider gibt es bezüglich der zu verwendeten Codecs unterschiedliche Meinungen
unter den Browserherstellern. Um die Möglichkeiten vom Audio Element vollständig
ausnutzen zu können wird angeraten MP3 und Ogg Vorbis als Codecs zu verwenden.

\subsubsection{Video}
Ein Video kann mittels folgendem Code auf einer Webseite angezeigt werden:
\begin{verbatim}
<video src="movie.ogg" width="640" height="360" controls>
<p>Your browser does not support the video element.</p>
</video>
\end{verbatim}
Auch beim Video Element ist die Frage nach dem Codec nicht einfach zu
beantworten. Das Video Element kann mit allen modernen Webbrowsern genutzt
werden, allerdings gibt es andere Probleme

\begin{itemize}
\item Apple Geräte können den Ogg Theora Codec aufgrund von Hardware Problemen
nicht wiedergeben
\item Opera und Firefox unterstützen den H.264 Codec aufgrund von Lizenz
Problemen nicht.
\end{itemize}

Erst vor kurzem brachte Google eine mögliche Lösung ist Spiel: Die Nutzung von
WebM. Dabei handelt es sich um einen lizenzfreie Videocodec. Firefox, Opera,
Chrome und auch IE haben bereits bestätigt, dass dieser Codec in Zukunft
unterstützt wird.
Einige der bekanntesten Video Streaming Anbieter sind bereits dabei HTML5 für
Videoinhalte zu nutzen, darunter unter anderem YouTube und Vimeo.

In Kombination mit einem Canvas Element können die Inhalte des Video Elements
auf verschiedenste Art manipuliert werden. Unter anderem könnte das Bild in
mehere Teile zerschnitten, Explosionen eingefügt und Filter angewendet werden.
Allerdings, benötigen komplexe Effekte und Manipulationen eine
Hardwarebeschleunigung. Flash nutzt zur effizienten Darstellung von Animationen
und Videos die Grafikkarte des Anwenders. Die Möglichkeiten von HTML5 sind noch
nicht soweit entwickelt und bietet keine native Möglichkeit um die Hardware des
Anwenders zu nutzen.

Um auch die Video Element effizient nutzen zu können, sollten die Quelldateien
mit mindestens zwei verschiedene Codecs zur Verfügung gestellt werden.

\section{Vor- und Nachteile}
Jeffrey Zeldman war schon im Jahr 2010 der Meinung, dass HTML, CSS und
JavaScript in Zukunft die Entwicklung von Rich Media Applikationen vorantreiben
wird. Einige der Vorteile werden hier aufgezählt.

\subsection{Plattformübergreifende Web-Technologien}
Web-Entwickler mussten vor der Umsetzung einer Webanwendung oder Webseite
entscheiden für welche Webbrowser entwickelt werden soll. Dieser Umstand
entstand aufgrund der ungeeigneten Konzeption früherer HTML und CSS Versionen.
Mit den neuen Webstandards (HTML5, CSS3 und JavaScript) ist es nun möglich
plattformübergreifende Anwendungen zu erstellen, die nicht nur auf allen
Desktopbrowsern sonder auch auf mobilen Browsern genutzt werden können.

\subsection{Optimiertes Mobile-Web}
Bei der Entwicklung von Anwendungen für mobile Geräte spielt die begrenzte
Bandbreite des Internets eine entscheidente Rolle. Mittels CSS3 können viele
Bilder durch Gradienten, Schatten und andere Effekte ersetzt werden.
Überdimensionierte Hintergrund- und Füllbilder werden somit überflüssig und
führen zu einem geringerem Verbrauch der zur Verfügung stehenden Bandbreite.
Einige HTML5 Features wie Local oder Session Storage erlauben es Daten
clientseitig zu speichern und Applikationen auch ohne Internet verfügbar zu
machen. Zusätzlich müssen Daten die bereits runtergeladen wurden nicht noch
einmal gealden werden.

\subsection{Rich Media ohne Plugins}
Die Audio, Video und Canvas APIs von HTML5 erlauben es Rich Media ohne
Verwendung von Plugins im Internet zu nutzen. Propritäte Technologien die in
diesem Bereich hauptsächlich verwendet werden haben mit HTML5 einen ernst zu
nehmenden Konkurrenten bekommen.

\subsection{Rich Media Applikationen}
Der Großteil der neuen HTML5 Features sind speziell für die Erstellung von Rich
Media Applikationen vorgesehen.

% =HTML5 Syntax
\section{Die Syntax von HTML5}

% =HTML5 Rich Media
\section{Rich Media}

% =HTML5 Web Forms 2.0
\section{Web Forms 2.0}

% =HTML5 Semantik
\section{Semantik}

% =HTML5 Schon Heute
\section{HTML5 schon heute}
%% =MOBILE
\chapter{Mobile Geräte}

\section{Die Entwicklung des mobilen Internets}
Die Optimierung von Applikationen und Webseiten für mobile Endgeräte gewinnt
immer mehr an Wichtigkeit. Wollte ein Nutzer vor der aktuellen Smartphonegeneration 
(iPhone und Android) unkompliziert mobil im Internet surfen oder E-Mails schreiben 
wollte, kam er nicht um ein Blackberry von RIM herum. RIM legte bei ihren Telefon
den Hauptfokus auf die E-Mail Funktion. Das Surfen auf einem mobilen Telefon
war meist durch zu kleine Bildschirme, langsame Verbindungsgeschwindigkeiten,
umständlicher Bedienung, zu hohen Kosten und Einschränkungen in der Darstellung
und Nutzung von Webinhalten stark eingeschränkt. Die Einführung der 
GSM-Erweiterungen GPRS und EDGE sowie die noch schnelleren, auf dem 
Mobilfunkstandard basierenden UMTS-Datenübertragungsverfahren HSDPA und
HSUPA sorgte, insbesondere mit dem Erscheinen des iPhone von Apple, für einen
langsamen Umschwung in der mobilen Internetnutzung. Das Bankhaus Morgan
Stanley hat im April 2010 eine 87-seitige Präsentation zum Thema "`Internet Trends"'
veröffentlicht, welche zum einen dem mobilen Internet ein schnelleres Wachstum 
gegenüber dem desktop-basierten Internet verspricht, und zum anderen, dass im
Jahr 2014 dieses in der Nutzung überholen wird.

\section{iOS von Apple}
\section{Android von Google}
%# HTML5
\chapter{HTML5}

\section{Was ist HTML5?}

%## HTML5 Geschichte
\section{Die Geschichte von HTML}
HTML, ``Hypertext Markup Language'', ist die allgegenwärtige
Auszeichnungssprache des World Wide Webs. Durch den geschickten Einsatz der
paar Elemente (fortan: Tags, Wörter umgeben von eckigen Klammern z.B. <html>)
die die Sprache unterstützt, konnten erstaunlich viele unterschiedliche
Netzwerke von verlinkten Dokumente erstellt werden. Von bekannten Seiten wie
Amazon, Ebay und Wikipedia bis hin zu personalisierten Blogs oder Webseiten,
die sich zum Beispiel auf Katzenbilder spezialisiert haben. HTML5 ist die
aktuellste Version dieser Auszeichnungssprache. Obwohl diese Version die
bisher umfangreichsten Änderungen mit sich bringt, ist es nicht die erste
Aktualisierung von HTML.
\newline\newline
Sir Tim Berners-Lee zeichnet sich für die Entwicklung von HTML und damit dem
Beginn des Internets verantwortlich. 1991 veröffentlichte er ein Dokument mit
dem Titel ``HTML Tags'', in dem er weniger als zwei Dutzend Elemente, die für
das Schreiben von Webseiten genutzt werden konnten, vorschlug.
Darunter auch einige Tags, die bis dato verwendet werden, wie zum Beispiel:
%
\begin{itemize}
  \item{Listen (<ul>, <ol>, <li>)}
  \item{Überschriften (<h1>, <h2>, ...)}
  \item{Paragraphen (<p>)}
  \item{Titel (<title>)}
  \item{Links (<a>)}
\end{itemize}
%
Die Verwendung von Tags war nicht Berners-Lees eigene Errungenschaft. Tags
wurden bereits in der SGML (Standard Generalized Markup Language) verwendet.
Anstatt einen komplett neuen Standard zu erfinden, erkannte Berners-Lee die
Vorteile, bereits existierende Standards weiterzuentwickeln - ein Trend der
auch in der Entwicklung der neuen HTML5 Spezifikation erkennbar ist.

\subsection{Von der IETF zur W3C: Der Weg zu HTML4}
Die erste offizielle Spezifikation war HTML 2.0, veröffentlicht durch die IETF,
die ``Internet Engineering Task Force''. Viele der neuen Features dieser
Spezifikation basierten dabei auf bereits veröffentlichte Implementierungen.
So bot der im Jahr 1994 marktführende Web Browser Mosaic Autoren von
Webseiten bereits die Möglichkeit, Bilder in Dokumente mittels eines <img>-Tags
einzubinden. Das <img>-Tag wurde in die HTML 2.0 Spezifikation inkludiert.
Die IETF wurde allmählich durch die W3C, dem ``World Wide Web Consortium'',
ersetzt. Folgende Aktualisierungen des HTML Standards wurden auf \url{
http://www.w3.org} veröffentlicht. In der zweiten Hälfte der Neunziger wurde
der Standard mehrmals überarbeitet bis 1999 HTML 4.01 veröffentlicht wurde.

\subsection{XHTML 1: Die Vermischung von HTML und XML}
Das nächste Update von HTML 4.01 trug den Namen XHTML 1.0. Das X steht dabei
für ``eXtensible''. Die Spezifikation von XHTML 1.0 war ident zu der von
HTML 4.01, somit wurden keine neuen Elemente und Attribute eingeführt. Die
Spezifikationen unterschieden sich nur durch die zu verwendende Syntax zum
Schreiben von Dokumenten. Während Autoren von HTML Dokumenten kaum
Einschränkungen im Schreibstil von Elementen und Attributen hatten, setzte
XHTML es voraus, dass die Regeln von XML, ``Extensible Markup Language'',
einer strikteren Auszeichnungssprache auf der viele Technologien des W3C
basierten, eingehalten werden. Aufgrund der strikteren Regeln einigten sich
Autoren auf einen gängigen Schreibstil, der auch unter dem HTML 4.01 Standard
Verwendung fand. Während Tags früher in Kleinbuchstaben, Großbuchstaben oder
einer Mischung aus beidem geschrieben werden konnten, verlangte ein valides
XHTML Dokument, dass alle Elemente und Attribute klein geschrieben werden.
Während XHTML 1.0 noch auf HTML basierte und lediglich die strikteren Regeln
von XML verwendete, war die XHTML 1.1 Spezifikation reines XML. Dieser Umstand
führte zu schwerwiegenden Problemen. XHTML 1.1 Dokumente konnten nicht mehr
unter dem Mime-Type \textit{text/html} definiert werden. Der bis dato
bekannteste Webbrowser Internet Explorer konnte jedoch Dokumente die mit einem
XML Mime-Type publiziert wurden, nicht darstellen.

\subsection{XHTML 2}
Das W3C war mit der vierten Version von HTML der Meinung, dass der HTML
basierte Ansatz seinen Zenith erreicht hat und setzte für die zukünftige
Version vollständig auf XML. Trotz der fast identen Namen von XHTML 1 und
XHTML 2, konnten die Unterschiede zwischen den beiden Spezifikationen nicht
größer sein. Anders als XHTML 1, war XHTML 2 nicht abwärtskompatibel mit
bereits existierenden Webinhalten oder vorangegangenen HTML Versionen. Es
sollte ein vollkommen neuer Standard werden und es zeigte sich, dass dieser
Weg nicht mit Erfolg gekrönt sein würde.

\subsection{Die Spaltung: WHATWG}
Innerhalb des W3C bildete sich eine Gruppierung die gegen die Ansätze
rebellierten. Namenhafte Vertreter von Opera, Apple und Mozilla waren mit der
eingeschlagenen Entwicklungrichtung des W3Cs nicht zufrieden. Ihnen war es
wichtiger, mehr Aufmerksamkeit auf Formate zur Entwicklung von Webapplikationen
zu richten. Bei einem Workshop in 2004 schlug Ian Hickson, der zu der Zeit bei
Opera arbeitete, eine Weiterentwicklung von HTML vor, mit der es möglich sein
soll, Anwendungen zu entwickeln. Der Vorschlag wurde durch die W3C abgelehnt.
Unzufrieden mit der Entscheidung bildeten die Rebellen eine eigene Gruppe:
Die ``Web Hypertext Application Technology Working Group'', oder kurz WHATWG.

\subsection{Der Weg von Web Apps 1.0 zu HTML5}
Von Beginn an operierte die WHATWG anders als das W3C. Während das W3C einen
kosensorientierten Ansatz anwendet, Themen werden vorgetragen, es wird
diskutiert und abgestimmt, wird auch bei der WHATWG diskutiert und abgestimmt,
allerdings liegt die finale Entscheidung, was in die Spezifikation komm und
was nicht, beim Editor. Der Editor ist Ian Hickson. Der W3C Ansatz klingt
demokratisch und fair, allerdings führt dieser auch dazu, dass der Prozess
stark verlangsamt wird. Bei der WHATWG hat jeder die Möglichkeit mitzuwirken,
da die letzte Entscheidung aber beim Editor liegt, entwickelt sich alles
schneller. Schon zu Beginn der Gründung von WHATWG wurden die Projekte in zwei
große Spezifikationen aufgeteilt: Web Forms 2.0 und Web Apps 1.0. Beide
Spezifikationen sollen die bisherige HTML Spezifikation erweitern. Mit
fortschreitender Entwicklung wurden alle Spezifikationen in eine einzige
implementiert und in HTML5 umbenannt.

\subsection{Die Wiedervereinigung}
Während HTML5 von der WHATWG weiter entwickelt wurde, setzte das W3C die
Entwicklung der XHTML 2 Spezifikation fort. Erst im Oktober 2006, schrieb Sir
Tim Berners-Lee einen Blog-Post, in dem er zugab, dass der Versuch, das
Internet von HTML auf XML zu übertragen, nicht funktionieren wird.
Nur wenige Monate später entschied sich das W3C eine neue HTML Working Group
zu bilden. Anstatt komplett von Anfang an zu beginnen, nutzen sie
glücklicherweise die bisherige Arbeit der WHATWG als Basis für zukünftige HTML
Spezifikationen. Dieser Umstand führte zu einer schwer überschaubaren
Situation. Das W3C arbeitete parallel an den beiden unterschiedlichen, nicht
kompatiblen Spezifikationen: XHTML 2 und HTML 5 (mit Leerzeichen).
Währenddessen arbeitete die WHATWG, an der Spezifikation für HTML5 (ohne
Leerzeichen) auf der die Arbeit der W3C aufbaut.

\subsection{HTML5: 2012 und 2022}
Heute ist der aktuelle Stand der HTML5 Spezifikation nicht mehr so
undurchsichtig wie er früher war, allerdings gibt es noch immer offene Fragen.
\newline\newline
Es gibt noch immer zwei Gruppen die an HTML5 arbeiten. Die wohl wichtigste
Frage für Webentwickler ist: Wann können wir es nutzen? In einem Interview gab
Ian Hickson an, dass die HTML5 Spezifikation frühestens 2022 den ``proposed
recommendation'' Status erreichen wird. Klingt nach einer langen Wartezeit,
allerdings bedeutet ``proposed recommendation'', dass die HTML5 Spezifikation
zwei mal komplett implementiert werden muss. Als Vergleich: HTML 4 existiert
nun seit über einem Jahrzehnt und hat noch nicht die gesetzten Features
erreicht. Wenn man nun den Umfang der Spezifikation betrachtet, klingt dieses
Datum hoch gesteckt. Browserhersteller sind nach wie vor nicht dafür bekannt,
dass existierende Standards so schnell wie möglich implementiert werden. Der
Internet Explorer benötigte mehr als ein Jahrzehnt um das \textit{abbr}-
Element richtig darstellen zu können.
\newline\newline
Für Webentwickler war das Jahr 2012 wesentlich wichtiger. 2012 erreichte die
HTML5 Spezifikation den ``candidate recommendation'' Status, der
gleichbedeutend ist mit "`fertig und abgeschlossen"'. Allerdings reicht das
alleine leider auch nicht aus. Wirklich entscheidend ist die Frage: Ab wann
unterstützen Webbrowser den neuen Standard? Schon die Veröffentlichung des
CSS 2.1 Standards zeigte, dass man nicht auf die Fertigstellung der
Spezifikationen warten sollte, sondern, wenn möglich, die Features nutzen
sollte sobald es möglich ist. Das selbe gilt auch für HTML5. Sobald Webbrowser
bestimmte Features der Spezifikation unterstützen, können diese auch jederzeit
verwendet werden. Man darf nicht vergessen, dass HTML5 keine komplett neu
entwickelte Sprache ist. Im Sinne der HTML Spezifikation ist es eher eine
Evolution als eine Revolution. Da HTML5 auf den früheren Versionen aufsetzt
und immer irgendeine Version des HTML Standards zur Erstellung von Webseiten
genutzt wird, wird bereits HTML5 genutzt.

\section{Features}
Mit HTML5 wird Entwicklern eine Fülle an neuen Funktionen und Features
angeboten. An dieser Stelle soll ein Überblick über die Funktionalitäten
gegeben werden, die diese Technologie auszeichnet.

\subsection{Grafiken und Animation}
\begin{itemize}
  \item{2D Zeichentools}
  \item{Animationstools}
  \item{3D Unterstützung}
  \item{3D Transformationen}
  \item{SVG (Scaleable Vector Graphics)}
  \item{Skalierbare Inhalte}
  \item{Web-Fonts}
  \item{Filtereffekte (z.B. Weichzeichnen)}
  \item{Präsentationen}
\end{itemize}
Die Spezifikation von HTML5 enthält bereits umfangreiche Ansätze zum Erstellen
von Grafiken und Animationen. So können mit CSS3 und dem neuen Canvas-Element
einfache und auch komplexe Grafiken erstellt und, mit dem Einsatz von der von
dem Canvas-Element zur Verfügung gestellten Schnittstellen, über JavaScript
animiert werden. Dank CSS3 können auch erstmals ohne Abstriche Schriftarten
integriert und verwendet werden. Ein Feature das bis zu diesem Zeitpunkt nur
von Flash angeboten wurde. Mit der Integration des SVG-Formats bietet auch
HTML5 eine Möglichkeit um skalierbare Inhalte in Webseiten einzubinden.
Einfache Filtereffekte können nun unkompliziert mit CSS3-Regeln erstellt
werden. Für komplexe Effekte kann das Canvas-Element in Kombination mit
JavaScript verwendet werden.

\subsection{Schnittstellen (APIs)}
\begin{itemize}
  \item{Zugriff auf das Filesystem}
  \item{Sockets}
  \item{Geolocation}
  \item{Datenaustausch}
\end{itemize}
HTML5 bietet wie Adobe Flash eine Vielzahl an Schnittstellen, um mit anderen
Systemen zu interagieren. Beispielsweise kann auf das Filesystem des
Computers zugeriffen werden, um Daten zu speichern oder zu laden.
Mittlerweile gibt es eine Unmenge an JavaScript Frameworks aus denen ein
Entwickler zur Umsetzung von Projekten zurückgreifen kann. Das wohl
bekannteste Framework ist jQuery. Frameworks helfen Entwicklern sonst schwer
zu realisierende Features einfacher umzusetzen, so bietet jQuery mehrere
Schnittstellen mit der ein Zugriff auf verschiedene Hardwarekomponenten, wie
zum Beispiel die zur Standordbestimmung, möglich ist. Schon frühere Versionen
von HTML unterstützten den Austausch von Daten über verschiedene Formate
wie XML, JSON oder andere.

\subsection{Multimedia}
\begin{itemize}
  \item{Videounterstützung}
  \item{Audiounterstützung}
  \item{Spracheingaben}
\end{itemize}
Eine der größten Neuerungen und Grund für den Schlagabtausch zwischen HTML5
und Adobe Flash bildet die Integration von nativen Audio- und Videoelementen
und deren Unterstützung. Beide Elemente können dabei über seperate
Schnittstellen angesprochen und gesteuert werden. Damit bietet sich eine
Möglichkeit an Multimedia-Inhalte in Webseiten einzubinden, ohne auf den Adobe
Flash Player zurückgreifen zu müssen. Einziger Wehrmutstropfen ist die noch
zum Teil unvollständige Unterstützung durch bestehende Browser, jedoch sollten
zukünftige Aktualisierungen dieses Problem beseitigen. In Kombination mit
weiteren Neuerungen von HTML5 und JavaScript lassen sich interessante und
komplexe multimediale Anwendungen erstellen. So kann Aufgrund der neuen
Spracheingabeschnittstelle auf das Mikrofon des Computers zugegriffen und
Spracheingaben verarbeitet werden.

\subsection{Sonstiges}
\begin{itemize}
  \item{
    Kompatibilität und Unterstützung auf unterschiedlichen Plattformen und
    Endgeräten
  }
  \item{Ansätze einer Programmiersprache}
\end{itemize}
HTML5 Inhalte lassen sich auf jeder und für jede Plattform erstellen.
Die Kompatibilität herzustellen liegt dabei nicht in der Hand des Entwicklers,
sondern in denen der unterschiedlichen Browserhersteller. Durch die steigende
Anzahl an Schnittstellen und den vielen Neuerungen nimmt JavaScript immer mehr
die Form einer komplexen Programmiersprache an, mit deren Hilfe komplexe
Anwendungen realisiert werden können.

\subsection{HTML5 Spezifisch}
\begin{itemize}
  \item{Semantisches Markup}
  \item{Suchmaschinenoptimierung}
\end{itemize}
Semantisches Markup ist wohl das Feature das HTML im allgemeinen auszeichnet.
Dank der angewandten Semantik ist es möglich, Inhalte mit, für den Computer
verständliche, Bedeutung zu versehen. Ein Feature das sich bereits im
Bereich der Suchmaschinenoptimierung als essenziel erwiesen hat und mit der
steigenden Bedeutung von Suchmaschinen nicht zu unterschätzen ist. Einige
interessante Neuerungen von HTML5 umfassen die Erweiterungen der bestehenden
Tags mit noch aussagekräftigeren Elementen.

\section{Aktueller Einsatz}
Keine andere Technologie konnte bis jetzt für einen ähnlich großen Hype sorgen
wie HTML5. Von vielen wird die neue Spezifikation als die einzige Technologie
angesehen, die notwendig sein wird, um in Zunkunft professionelle,
anspruchsvolle und interaktive Webinhalte zu erstellen. Gegenwärtig werden
diese Erwartungen allerdings noch nicht vollständig erfüllt. Die aktuelle
Version, die veröffentlicht wurde, ist lediglich ein Entwurf für den finalen
Webstandard. Bis dieser Standard verabschiedet wird, dauert es noch bis ins
Jahr 2022. Das hat zur Folge, dass Entwickler mit Änderungen im
Funktionsumfang und der bestehenden Features rechnen können.
\newline\newline
Nichts desto trotz kann und soll der neue Standard für die Entwicklung
von Webseiten und -applikationen genutzt werden. Mit HTML5 wurden viele
Altlasten älterer HTML Versionen beseitigt und einige neue und zukunftsichere
Funktionen hinzugefügt. In einigen Bereichen wie Multimedia, Interaktivität
und Schnittstellen steht HTML5 dem bisherigen Platzhirsch Adobe Flash in
den erwähnten Bereichen kaum noch nach. Im direkten Vergleich zu Adobe Flash
löst HTML5 einige Aufgaben deutlich besser und verbraucht dabei auch noch
weniger Ressourcen und Bandbreite.
\newline\newline
Mit der kürzlichen Umstellung großer Videoportalen (beispielsweise Youtube)
von Adobe Flash auf HTML5 lässt sich eine eindeutige Bewegung feststellen.
Allerdings wird bis jetzt die HTML5 Variante lediglich als zusätzliche
Möglichkeit neben Flash angeboten. In Zukunft soll jedoch komplett auf die
Verwendung von Flash verzichtet werden. Was die Umstellung bei Webseiten
betrifft, geht diese wesentlich gemächlicher voran.
\newline\newline
Im Gegensatz zu Adobe Flash ist für die Verbreitung von HTML5 nicht nur
der Einsatz des Standards sondern auch die Unterstützung der Features durch
die verschiedenen Browser notwendig. Leider zeigt sich hierbei eine der
Schwächen von HTML5: Längst nicht alle Features werden von allen Browsern
unterstützt. Ganz im Gegenteil, interessante neue Funktionen werden
von vielen Browsern nur teilweise oder garnicht unterstützt.
Das führt unweigerlich zu sehr unterschiedlichen Darstellungen einer
Webseite auf mehreren Browsern und damit zu einem Mehraufwand und höheren
Kosten in deren Entwicklung und Testung.

\section{Technologieverbund}
Wenn von HTML5 gesprochen wird, ist in den häufigsten Fällen der Verbund
aus der Markupsprache HTML in der Version 5, der deklarativen Sprache
CSS in der Version 3 und der Skriptsprache JavaScript (auch bekannt
als ECMAScript) gemeint. Alle drei Technologien greifen ineinander und
bieten in ihrer Gesamtheit die neuen Features von HTML5.
% Bild mit Triangle HTML Struktur/Inhalt, CSS Präsentation/Darstellung, JS Logik

\subsection{HTML5}
HTML ermöglicht es, Inhalte im Web zu strukturieren und darzustellen.
Zusätzlich bietet HTML Funktionen um Texte, Bilder, Hyperlinks,
Videos, Audio und andere multimediale Inhalte semantisch auszuzeichnen
und damit für andere Systeme, zum Bespiel dem Computer, verständlich
zu machen. Jegliche Inhalte, die im Web dargestellt werden sollen, müssen
in die HTML-Struktur eingebunden werden, selbst solche, die zur Darstellung
Plug-Ins von anderen Anbietern wie beispielsweise Adobe Flash benötigen.
\newline\newline
Die gewünschten Inhalte werden dabei in verschiedene HTML-Elemente
eingebettet. Jedes dieser HTML-Elemente erfüllt einen bestimmten Zweck und
dieser reicht von Überschriften und Paragraphen bis hin zu Hyperlinks und
Bildern. Zusätzlich können diese Elemente durch verschiedene Attribute noch
genauer ausgezeichnet werden. Generische Elemente wie das div-Element,
ergänzen den begrenzten Elemente-Umfang, um nicht eindeutig zuordenbare
Inhalte auf einer Webseite zu platzieren. Über diesen Weg kann jede
Form von Inhalt erstellt werden.
\newline\newline
Für HTML5 wurden verschiedene Ziele festgelegt die es zu erreichen galt.
Diese Ziele dienten zusätzlich dem Zweck, den neuen Standard konkurrenzfähig
und zukunftssicher zu machen. Die Ziele waren:
%http://www.w3.org/TR/html-design-principles/
\begin{description}
  \item[Kompatibilität:]
  Bereits bestehende Inhalte müssen weiterhin verwendet werden können. Neue
  Funktionen des Standards dürfen die bestehenden Inhalte nicht verändern oder
  beeinflussen.
  \item[Verwendbarkeit:]
  Neue Funktionen sollen bestehende Probleme lösen und die Entwicklung
  einfacher machen.
  \item[Sicherheit:]
  Ein wichtiger Aspekt, der im Web immer mehr an Bedeutung gewinnt.
  \item[Konsistenz:]
  Sowohl in HTML als auch in XHTML sollen XML-Elemente genutzt werden können.
  Die Sprachen greifen dabei auf eine gemeinsame Dokumentation zurück.
  \item[Vereinfachung:]
  HTML soll deutlich vereinfacht und durch konkrete Fehlermeldungen schneller
  und einfacher korrigierbar werden.
  \item[Universalität:]
  Inhalte sollen in jeder Sprache und auf allen Endgeräten sowie Systemen
  darstellbar sein.
  \item[Zugänglichkeit:]
  Inhalte und Funktionen sollen jedem zugänglich sein. Dies entspricht dem
  OpenSource-Gedanken.
\end{description}
%
Bis dato wurden bereits einige Ziele erreicht. Bis zur offiziellen
Verabschiedung des Standards im Jahr 2022 werden mit Sicherheit alle Ziele
erreicht werden.

\subsection{CSS3}
Mittels ``Cascading Stylesheets'', oder kurz CSS, wird die Darstellung von
Elementen in einem HTML Dokument kontrolliert. So ist es möglich die
Positionierung und Attribute wie Farbe, Schattierung oder ähnliches, von
Elementen exakt zu definieren. CSS ist damit das Werkzeug der Wahl um
Designvorgaben möglichst genau umzusetzen.
\newline\newline
Die Trennung von HTML und CSS symbolisiert die Trennung von Inhalt und
Design. Dies ermöglicht eine unterschiedliche Darstellung der selben Inhalte
durch unterschiedliche Stylesheets für verschiedene Endgeräte oder den Druck,
ohne dabei den Inhalt selbst anpassen zu müssen.
\newline\newline
Der Name Cascading Stylesheets ist abgeleitet von den zur Verfügung stehenden
Vererbungsregeln, die das Format zulässt. Durch die Vererbungsregeln ist es
möglich, ein große Auswahl an Elementen gleich aussehen zu lassen und
dennoch jedes beliebige Element individuell anzupassen. Die von
HTML gelieferten Attribute ID und Class dienen zur Unterscheidung von
einmalig auftretenden Elementen (id) und ganzen Gruppierungen (class).
\newline\newline
Mit dem Level 3 des Cascading Stylesheet Standards kommen viele neue
Eigenschaften hinzu, die die Gestaltung von Elementen noch schneller und
einfacher machen. Mit CSS2 konnten vorwiegend einfache Darstellungsparameter
wie Farbe, Rahmen, Innenabstand und Aussenabstand sowie die genau
Positionierung im Dokumentenfluss definiert werden. CSS3 ermöglicht nun
etliche grafische Effekte wie Verläufe, Schatten oder abgerundete Ecken.
Zusätzlich zählen zu den wichtigsten neuen Funktionen von CSS3 die Erstellung
von einfachen Animationen, das Einbinden und Modifizieren von Schriftarten
sowie Transformationen und Skalierungen von Inhalten.
\newline\newline
Seit dem Jahr 2000 ist der CSS3 Standard in Entwicklung und die komplette
Unterstützung durch Browserhersteller schreitet immer weiter voran. Vor
allem durch das Aufkommen von HTML5 werden die Eigenschaften immer stärker in
den Mittelpunkt gerückt. Aus diesem Grund werden die Funktionen von CSS3 oft
dem HTML5 Standard zugeordnet.

\subsection{JavaScript}
Die Skriptsprache JavaScript übernimmt in dem Technologieverbund die logische
Ebene. Die Verwendung von HTML als eigenständige Technologie bietet nur die
statische Erstellung und Darstellung von Inhalten. Durch JavaScript lassen
sich HTML-Elemente allerdings dynamisch verändern. Ein Trend der sich bei
der Entwicklung von modernen Webseiten durchgesetzt hat und mittlerweile
nicht mehr weg zu denken ist. JavaScript wurde 1995 offiziell von Brendan Eich
eingeführt und ist aktuell in der Version 1.8.5 verfügbar. Allerdings
unterstützen ein Großteil lediglich die Version 1.5 vollständig.
\newline\newline
Die Skriptsprache ermöglicht die Manipulation von HTML-Elementen auf der
im Browser aufgerufenen Seite, das dynamische Erstellen neuer Elemente oder
das dynamische Nachladen von Inhalten mit Hilfe von XML und AJAX. Besonders
für einfache Effekte wie ein Dropdown-Menü oder das Einblenden von Tooltips
bietet sich der Einsatz von JavaScript an.
\newline\newline
Als dritte Säule der Webseitenerstellung kümmert sich JavaScript daher vor
allem um die Dynamik und Interaktivität von modernen Webseiten. Diese
Funktionalitäten werden heutzutage von Nutzern erwartet. Die steigenden
Erwartungen und die immer komplexer werdenden Webseiten fordern auch, dass
JavaScript diesen gewachsen ist. Diese Anforderungen und die fehlende
Unterstützung der aktuellsten Skript-Versionen machen die Entwicklung
unverhältnismäßig schwer, saubere Lösungen für an sich einfache Probleme zu
entwickeln. Für diesen Fall kommen die sogenannten JavaScript Frameworks ins
Spiel. Frameworks bieten Anhand von kompatibler JavaScript Versionen neue
Funktionalitäten. Besonders um einfach, sicher und effizient Animationen oder
Effekte zu realisieren oder Elemente zu manipulieren bietet sich die Nutzung
von Frameworks an. Einige der bekanntesten Frameworks sind dabei Prototype,
MooTools und das mit Abstand am häufigsten verwendete Framework jQuery.
\newline\newline
Nur durch das Zusammenspiel des Grundgerüsts HTML, des Designs CSS3 und der
Interaktivität durch JavaScript wird HTML5 zu dem Gesamtprodukt, welches in
Zukunft die Entwicklung des Webs prägen wird und zu einem ernst zu nehmenden
Konkurrenten für bereits etablierte Softwarelösungen wie Adobe Flash werden
lässt.

\section{Notwendigkeit von HTML5}
Die letzte Version von HTML wurde gut 10 Jahre lang eingesetzt. Mit dem
unglaublichen schnellen Wachstum im Bereich der mobilen Endgeräte, wie
Smartphones und Tablets, und den steigenden Erwartungen der Nutzer im Bezug
auf Inhalte und Funktionalitäten konnte die bestehende Version den gegebenen
Ansprüchen nicht mehr gerecht werden. Um sowohl die geforderten Anforderungen
zu erfüllen und zukunftssicher zu sein ist eine neue Version der Sprache
notwendig. HTML5 kann diese Version werden.
\newline\newline
Die bisher vergleichweise rasante Entwicklung von HTML5 wird sich daher auch
in Zukunft kaum verlangsamen. Anders als bei einer proprietären Software
beteiligen sich viele der führenden Unternehmen im Bereich Internet an der
Entwicklung des neuen Standards. Damit kann man von einer umfassenden
Unterstützung durch verschiedene Browser, Endgeräte und Schnittstellen
ausgehen.
\newline\newline
Um technische Neuerungen und moderne Webseiten zu gestalten mussten bisher
immer die Funktionen von Adobe Flash oder ähnlicher Software genutzt werden.
Die neuen Features von HTML5 sollen hingegen eine Möglichkeit bieten
komplett auf Software-Lösungen von Drittanbieteren zu verzichten und dabei
einen einheitlichen Standard für alle Geräte und Plattformen zu schaffen.
\newline\newline
Vor allem im Bereich der Multimediainhalte wie Videos, Audio, Animationen, 3D
und Interaktion musste Adobe Flash genutzt werden. Zu beachten ist dennoch,
dass Software von Drittanbietern neben technischen Vorteilen auch einige
Nachteile mit sich bringen. Einer der größten Nachteile stellt die kaum
oder nur mit sehr viel Aufwand verbundenen Suchmaschinen-Optimierung von
Flash-Inhalten dar. Ein Problem, dass in einer Zeit in der Online-Marketing
immer mehr an Bedeutung gewinnt, nicht zu unterschätzen ist. Weiters
werden laufend neue Sicherheitsproblemem bei den angebotenen Plug-Ins, die
notwendig sind um entsprechende Inhalte darzustellen, bekannt und nur langsam
oder auch unzureichend behoben. Aufgrund der vielen Nachteile wurde der Ruf
der Kritiker, von Adobe Flash und proprietärer Software, nach einer
Alternative immer lauter.
\newline\newline
Mit der Entscheidung von Apple Adoble Flash auf ihren mobilen Endgeräten
nicht mehr zu unterstützen, führte außerdem dazu, dass schnell Alternativen
entwickelt werden mussten. HTML5 stellt dabei eine optimale Lösung dar.
%# HTML5 Canvas
\chapter{HTML5 Canvas}

% Überblick Canvas Element (Was ist es?)
% Wie funktioniert es (Vergleich Canvas/Flash)
% Überblick Canvas API (Was kann es?)
% Beispiele (Wie wird es gemacht?)
% Wo wird es schon verwendet?

\section{Einleitung}
Das Canvas Element bildet eines der größten Teile der HTML5 Spezifikation.
Es bietet eine API für das Zeichnen von 2D Objekten, wie Linien, Füllungen,
Bilder oder Text. Ein Vergleich mit MS Paint ist sehr gut um sich die
Möglichkeiten mit dem Canvas Element vorstellen zu können. Tatsächlich
gibt es bereits mehrere Implementationen, basierend auf dem Canvas Element,
die Applikationen wie MS Paint vollständig ersetzen können. Es geht bereits so
weit, dass laufend neue Zeichenapplikationen entstehen, die sich darum
bemühen, als vollwertige Vektor-Zeichenprogramm, ähnlich wie Adobe
Illustrator, angesehen zu werden. Das bemerkenswerte daran ist, dass es
auch einige Anwendungen gibt, die ebenfalls auf mobilen Endgeräten, wie
iPhones oder Android Smartphones, funktionstüchtig sind.

% Pics, Canvas MS Paint, Sketchpad, Harmony

\section{Geschichte}
Das Konzept des Canvas Elements wurde durch Apple eingeführt um in Mac OS X
Webkit Widgets für das Dashboard zu erstellen. Vor der Einführung des Canvas
Elements konnten Zeichen Schnittstellen im Browser nur dann genutzt werden,
wenn Plugins wie Flash von Adobe, Scaleable Vector Graphics (SVG), die
Vector Markup Language (VML), welche nur in Internet Explorer vorhanden war,
oder diverse JavaScript Bibliotheken eingesetzt wurden. % (Mary Nyamor S.25)
\newline
Das Canvas Element ist eine Umgebung für das Erstellen von dynamischen Bildern.
Das Element selbst ist dabei sehr einfach zu verwenden. Lediglich die
Dimensionen des Canvas und eine id, über die das Canvas anschließend über
JavaScript referenziert werden kann, müssen festgelegt werden.
%
\lstset{style=HTML5}
\begin{lstlisting}
  <canvas id="myCanvas" width="360" height="240">
  </canvas>
\end{lstlisting}
%
Alles was sich zwischen den sich öffnenden und schließenden <canvas> Tags
befindet, wird nur Browsern angezeigt, die das Canvas Element nicht
unterstützen.
%
\begin{lstlisting}
  <canvas id="myCanvas" width="360" height="240">
    <p>Canvas is not supported by your browser.</p>
  </canvas>
\end{lstlisting}
%
Um auf dem Canvas zeichnen zu können muss JavaScript verwendet werden. Dabei
muss zunächst das entsprechende Canvas über seine id ausgewählt und
anschließend der gewünschte Kontext definiert werden. Die Auswahl des Kontexts
beeinflusst die Funktionen der genutzen API:
%
\lstset{style=JavaScript}
\begin{lstlisting}
  var canvas = document.getElementById("myCanvas");
  var context = canvas.getContext("2d");
\end{lstlisting}
%
Beim Verfassen dieser Arbeit standen nur der 2D oder der WebGL (3D) Kontext zur
Verfügung. Mit der Auswahl des gewünschten Kontexts ist das Canvas Element
bereit um für das Zeichnen von Bilder genutzt zu werden. Die 2D API verfügt
über alle Funktionen die bei einem herkömmlichen, teils kostenpflichtigen,
Grafikprogramm wie zum Beispiel Adobe Illustrator erwartet werden: Es können
Linien, Konturen, gefüllte Flächen, Verläufe, Schatten, Formen und Bézier
Kurven gezeichnet werden. Der wesentliche Unterschied zwischen dem Canvas
Element und einem Grafikprogramm besteht allerdings darin, dass das Canvas
Element kein grafisches User Interface, zum Erstellen von Bildern, zur
Verfügung stellt, sondern jeder Pixel per JavaScript programmiert werden muss.

\subsection{Mit Code zeichnen}
Um die Farbe einer Kontur bzw. einer Linie zu definieren ist folgender Code
notwendig:
%
\begin{lstlisting}
  context.strokeStyle = "\#990000";
\end{lstlisting}
%
Alles was nun folgend auf dem Canvas gezeichnet wird, hätte eine rote Kontur.
Um ein Rechteck mit Kontur und ohne Füllung auf das Canvas zu Zeichnen gibt es
die strokeRect Methode und ist folgendermaßen definiert:
%
\begin{lstlisting}
  strokeRect(left, top, width, height);
\end{lstlisting}
%
Um nun ein rotes Rechteck zeichnen zu können, das 20 Pixel vom linken Rand und
30 Pixel vom oberen Rand des Canvas entfernt ist, 100 Pixel breit und 50 Pixel
hoch ist, ist folgender Code notwendig:
%
\begin{lstlisting}
  context.strokeRect(20, 30, 100, 50);
\end{lstlisting}
%
Dabei handelt es sich noch um ein sehr einfaches Beispiel um das Canvas zum
Zeichnen zu nutzen. Der 2D Kontext bietet eine sehr umfangreiche Auswahl an
vordefinierten Methoden wie fillStyle, zum Festlegen von Füllungen, fillRect,
um ein gefülltes Rechteck zu zeichnen, lineWidth, zum Festlegen der
Linienstärke, shadowColor, zum Festlegen der Schattenfarbe, und viele mehr.
Theoretisch betrachtet kann jedes Bild das mit einem Programm wie Adobe
Illustrator erstellt werden kann, ebenfalls mit dem Canvas Element realisiert
werden. Praktisch sieht das anders aus. Ein Versuch würde sich als
äußerst arbeitsaufwendig und in einer Unmenge an JavaScript äussern.
Abgesehen davon, wurde das Canvas Element nicht zur Reproduktion von
aufwendigen Pixel- oder Vektorgrafiken entwickelt.

\subsection{Canvas: Was bringts?}
Im vorherigen Abschnitt wurde ein Beispiel, die Reproduktion von Bildern,
genannt, wofür das Canvas Element nicht verwendet werden sollte. Die wirkliche
Macht des Canvas Elements liegt in der Möglichkeit, dass dessen Inhalt
jederzeit, aufgrund von zum Beipiel Nutzereingaben, verändert werden kann.
Gerade das Reagieren auf Nutzereingaben ermöglicht die Erstellung von
Anwendungen und Spiele die ursprünglich eine Technologie wie Adobe Flash
benötigten.
\newline
Eines der ersten Vorzeigebeispiele für die Möglichkeiten des Canvas Elements
wurde von Mozilla Labs entwickelt. Bespin: Ein Code Editor der innerhalb des
Browserfensters läuft. % (https://bespin.mozilla.com) + bild
Eine sehr eindrucksvolle und nützliche Applikation, allerdings von vielen
Entwicklern als perfektes Beispiel betrachtet, wofür das Canvas Element nicht
genutzt werden sollte.



% =FLASH
\chapter{Adobe Flash}

\section{Was ist Flash?}
Während die Popularität von HTML5 immer weiter steigt, werden immer mehr Vergleiche
mit "`Flash"' durchgeführt. Hierbei handelt es sich um ein Produkt der Firma Adobe, die 
unter anderem auch für Software wie Photoshop oder Illustrator verantwortlich ist. Allerdings
bezieht sich das Wort Flash in den meisten Fällen nicht auf die gleichnamige 
Entwicklungsumgebung der Firma, sondern der Flash Player, der das Abspielen von 
proprietären SWF-Dateien im Browser ermöglicht. Anfänglich war Flash ein
Programm um Illustrationen und einfache Animationen zu erstellen und wurde
1997 vonder Firma Macromedia, welche im Jahr 2005 von Adobe aufgekauft wurde,
entwickelt.

\section{Die Geschichte von Flash und ActionScript}
Seit der Veröffentlichung von Flash in 1996 wurde Flash und ActionScript gemeinsam
weiter entwickelt. Mittlerweile stellt die Kombination aus Design- und Animations-Tools 
in Flash und den interaktiven Möglichkeiten von ActionScript eine der mächtigsten, 
vielseitigsten und beliebtesten Entwicklungsumgebung dar. Allerdings war ActionScript 
zu beginn eine sehr bescheidene Programmiersprache.
\newline\newline
Die ersten drei Versionen von Flash boten noch keine Programmiertools an und 
interaktive Abläufe wurden durch Drag-and-Drop Optionen als Actions festgelegt. 
Lediglich die Navigation durch die Zeitleiste und das Erstellen von Links war 
über diese Actions möglich.
\newline\newline
Flash 4 war die erste Version, die die Möglichkeit bot Code mittels einer einfachen
Skriptsprache zu schreiben, welche inoffiziell bereits ActionScript genannt wurde.
Mit Flash 5 entwickelte sich ActionScript weiter und wurde zur offiziellen Skriptsprache.
Mit jeder folgenden Version von Flash wurden auch die Möglichkeiten von ActionScript erweitert.
ActionScript ermöglichte die interaktive Kontrolle von Text, Animation, Sound, Video, Data und
mehr. 2003 wurde ActionScript 2.0 vorgestellt, und dessen Möglichkeiten waren schon
vergleichbar mit objekt-orientierten Programmiersprachen wie Java oder C\#.
\newline\newline
Professionelle Programmierer interessierten sich immer mehr für ActionScript als 
Entwicklungstool, jedoch zeigte sich, dass trotz der vergleichbaren Möglichkeiten 
ActionScript in der Performance noch nicht mit seiner Konkurrenz mithalten konnte.
Der Grund dafür war, dass jede Version von ActionScript auf der vorherigen aufbaute.
Der Flash Player war ursprünglich nicht für aufwändige Anwendungen und komplexe Spiele
konzipiert, dennoch begannen Entwickler den Flash Player dafür zu nutzen.
Damit wurde es klar, dass eine neue Version von ActionScript von Beginn an neu entwickelt
werden musste um den Erwartungen gerecht zu werden.
\newline\newline
2006 wurde von Adobe ActionScript 3.0 vorgestellt, welches bedeutend mehr neue 
Funktionalitäten und eine starke Performancesteigerung mit sich brachte. Flash CS3 
war die erste Version von Flash die ActionScript 3.0 unterstützte. Flash CS4 
fügte neue Funktionen zu ActionScript 3.0 hinzu, unter anderem neue 3D Möglichkeiten,
Kontrollfunktionen für Animationen und ActionScript Klassen für Adobe AIR. Flash CS5 und
Flash CS6 erweiterten die Funktionalitäten unter anderem mit der Unterstützung von Controller 
und anderen Geräten, inklusive Multitouch-Funktionen und Touch-Screen Geräten.

\section{Flash Verfügbarkeit auf Computer}
Trotz propritärer Technologien hat Adobe es geschafft, ihren Flash Player auf 99\%
Prozent aller Systeme unterzubringen. Gründe dafür waren die Fähigkeiten von Flash
mit denen Ziele erreicht werden konnten, die mit Hilfe von offenen Standards nicht
möglich waren. In den Anfangstagen von Flash wurde es vor allem für Introanimationen,
welche auf Startseiten eingesetzt wurden, interaktive Navigationen, die auf die Aktionen
des Users mit z.B. Animationen reagierten, oder für die allseits bekannten Werbebanner.
Im Gegensatz zu Flash war die Erstellung ähnlicher Funktionen mit standardkonformen 
Technologien wie HTML 4.01, JavaScript oder CSS nur schwer oder überhaupt nicht 
realisierbar. Selbst heute werden für komplette Webseiten, Spiele und vor allem
Werbebanner mittels Flash realisiert. Ein weiterer Vorteil von Flash ist, dass es so einfach
zu erlernen und bedienen ist, das auch Neueinsteiger sehr schnell zu ansehnlichen 
Ergebnissen kommen.

\section{Flash Verfügbarkeit auf mobilen Geräten}

\section{Möglichkeiten}
\subsection{Spiele}
\subsection{Animationen}
\subsection{Entertainment}

\section{Die Einflüsse von Flash auf das Web}
\subsection{Die Wurzeln und Entwicklung von Flash}
\subsection{Stärken von Flash}
\subsection{Wichtiger Beiträger für die Entwicklung des Webs}
\subsection{Apples Argumente gegen Flash}
Ein oft verwendete Aussage und immer wiederkehrendes Streitthema vieler Diskussionen 
ist sicherlich, ob HTML5 langfristig Flash als Webentwicklungswerkzeug ablöst. Apple
spielt hierbei eine große Rolle. Das Unternehmen verzichtete seit der Einführung ihrer 
iDevices (iPhone, iPod touch, iPad) auf die Nutzung von Flash. Vor allem die fehlende Möglichkeit
Flash auf dem iPad zu verwenden führte vermehrt zu Kritik. Der User sei eingeschränkt und
ohne Flash handelt es sich nicht um das "`echte Internet"', welches Apple mit ihren Produkten
anpreist. Der verstorbene Steve Jobs, CEO von Apple, hat daraufhin einen offenen Brief 
verfasst, der den Nutzern erklärt, warum die Firma auf Flash verzichtet hat und auch
weiterhin darauf verzichten wird. Steve Jobs Hauptargumente waren und gelten noch bis
heute:
\begin{itemize}
	\item[1]{
		Flash ist ein 100 Prozent proprietäres System von Adobe, welches trotz der
		allgemeinen Verfügbarkeit des Flash Players die Zügel in der Hand hat.
		Adobe könnte ganz alleine über die Richtung entscheiden, in die Flash
		geht oder die Preisgestaltung ändern. Apple möchte das Internetstandards
		offen sein sollen, selbst wenn es auf die eigenen Geräte zutrifft.
	}
	\item[3]{
		Flashprodukte sind laut einer Analyse des Softwarehauses Symantec mit
		am anfälligsten für Sicherheitsprobleme. Sie seien einer der Hauptgründe
		für Abstürze auf dem Mac, auch die Performance auf Mobiltelefonen würde
		zu wünschen übrig lassen.
	}
	\item[4]{
		Flash beansprucht die Akkulaufzeit auf mobilen Geräten, da das Dekodieren
		von Videomaterial über die Software laufen muss. H.264 würde über die
		Hardware dekodiert werden.
	}
	\item[5]{
		Flash wurde nicht für Touchscreens geschaffen. Viele Webseiten müssten
		ihre Seite komplett neu konzipieren und entwickeln. Entwickler sollten
		aber dann doch auf modernere Technologien wie HTML5, CSS3 und
		Javascript zurückgreifen.
	}
\end{itemize}
Ein kommerzieller Hersteller wirf nun einem anderen kommerziellen Hersteller vor zu
proprietär zu sein - das mag durchaus merkwürdig erscheinen. Vor allem weil die 
Kunden von Apple auch im eigenen System eingezäunt und der Kontakt zu anderen 
Systemen über Schnittstellen so gering wie möglich gehalten wird un bei Adobe nur 
von der theoretischen, aber unwahrscheinlichen Möglichkeite ausgegangen wird, dass
diese die Vormachtstellung ihres Browserplugins missbrauchen könnten.
\newline\newline
Auch die Performanceprobleme scheinen bei Googles Smartphone Betriebssystem
Android seit der Version 2.2 gelöst worden zu sein. Denn seitdem läuft
Flash stabil und schnell, ohne den Akku zu sehr zu beanspruchen.
\newline\newline
Die Behauptung das Flash nicht für Touchscreen ausgelegt wirden ist, ist an und für
sich wahr, das gleiche gilt allerdings auch für HTML. Auch hier gibt es eben so viele
mausbasierende Aktionen, wie z.B. Rollover Animationen.
% =VERGLEICH

% Links
% http://venturebeat.com/2012/01/31/html5-versus-flash-infographic/
% http://www.periscopic.com/#/news/2011/05/our-research-into-flash-and-html5-which-one-is-right-for-your-project/
% http://www.actionscript.org/resources/blogs/53/A-good-comparison-between-HTML5-and-Flash.html
% http://blog.accusoft.com/posts/2012/october/html5-vs-flash-what-do-you-need-to-know-part-1.html
% http://blog.accusoft.com/posts/2013/january/html5-vs-flash-infographic.html
% http://blog.accusoft.com/posts/2012/october/html5-vs-flash-what-do-you-need-to-know-part-2.aspx
% http://readwrite.com/2010/03/09/does_html5_really_beat_flash_surprising_results_of_new_tests#awesm=~obelLBPH3UgwgC

\chapter{HTML5 Canvas und Adobe Flash: Vergleich und mögliche Auswirkungen}

\section{Vorraussetzungen}
Sowohl Nutzer als auch Entwickler benötigen um HTML5 und Adobe Flash einsetzen
zu können verschiedene Tools. Auf Entwicklerseite sind die Möglichkeiten sehr
vielfältig, sodass eine vollständige Auflistung den Rahmen dieser Arbeit
sprengen würde. Auf Nutzerseite ist die Auswahl der Tools ebenfalls breit
gefächert, allerdings ist die Anzahl der notwendigen Programme vergleichweise
geringer und kann daher genauer untersucht werden.
\newline\newline
Die Entwicklung von Flash Anwendungen setzt einen Editor vorraus, der in der
Lage ist ActionScript Code zu kompilieren. Dabei gibt es eine umfangreiche
Auswahl an verschiedenen Editoren, die nicht zwingend von Adobe kommen müssen.
Adobe bietet mit Adobe Flash CS6 eine Entwicklungsumgebung, welche sämtliche
Funktionen von Flash und ActionScript unterstützt und dem Entwickler neben dem
Editieren des ActionScript Codes eine grafische Oberfläche bietet, die es
vereinfacht, Animationen zu erstellen oder Grafiken einzubinden, allerdings
auch sehr kostspielig ist. Über die kostenlos erhältlichen Editoren können
zwar alle Funktionen von ActionScript genutzt werden, aber es werden keine
Tools zur Erstellung von Grafiken und Animationen angeboten.
\newline\newline
Im Vergleich dazu benötigen HTML5 Entwickler lediglich einen Editor, mit dem
textbasierte Inhalte erstellt werden können und einen Browser, der diese
Inhalte korrekt darstellt. Die Lösungen hierfür sind nahezu grenzenlos,
weshalb Entwickler aufgrund ihrer Präferenzen die für sie beste Lösung wählen
können. Der komplette Aufbau der Webseite wird im Editor in der entsprechenden
HTML5 Syntax geschrieben. Multimediale Inhalte lassen sich ebenfalls über das
Markup einbinden, und können zuvor mit den entsprechenden Programmen den
Wünschen angepasst werden. Adobe ist dabei eines der ersten Unternehmen, die
Softwarelösungen bieten, mit denen Webseiten und Animationen in HTML5 mittels
einer grafischen Oberfläche erstellt und bearbeitet werden können. In Zukunft
ist es möglich, dass andere Unternehmen diese Idee aufgreifen und ähnliche
Programme anbieten werden, sodass HTML5 im Bezug auf grafische
Entwicklungsumgebungen Adobe Flash in nichts nachstehen wird.
\newline\newline
Die technischen Vorraussetzungen für Nutzer von Flash-Inhalten sind relativ
einfach zu erfüllen. Der Nutzer müssen lediglich eine aktuelle Version des
Adobe Flash Players auf seinem System installiert haben, sei es als Browser
Plug-In oder als Standalone Version, und das verwendete System muss die
Hardware-Mindestanforderungen von Adobe erfüllen. Da diese
Mindestanforderungen äusserst gering sind, sollten alle aktuellen Systeme,
unabhängig ob mobil oder Desktop, in der Lage sein Flash-Inhalte abzuspielen.
Da der Adobe Flash Player auf jedem gängigen System verwendet werden kann,
ist auch die Entwicklung über mehrere Plattformen in der Regel kein Problem.
\newline\newline
Für die Nutzung von HTML5 Inhalten wird lediglich ein Browser vorrausgesetzt,
der diese korrekt darstellen kann. Der Entwickler hat dafür zu
sorgen, dass jeder Browser die Seite soweit korrekt darstellt, sodass der
Nutzer auf die Inhalte zugreifen kann. Das Problem hierbei ist, dass der
Browsermarkt sehr stark fragmentiert ist und diese sowohl von Hersteller zu
Hersteller als auch von Version zu Version eine unterschiedlich
fortgeschrittene Unterstützung von HTML5 anbietet. Die Hardware-
Mindestanforderungen sind dabei an den jeweiligen Browser gebunden. Allerdings
sind auch diese sehr gering wodurch auch schon auf mobilen Systemen Browser
integriert sind die eine sehr fortgeschrittene Unterstützung von HTML5
gewährleisten.

\section{Betriebssysteme}
HTML5 ist ein offener Standard. Das bedeutet, dass die Technologie nicht
von einem großen Softwarehersteller entwickelt wird und somit für die
Öffentlichkeit frei zugänglich ist. Das hat zur Folge, dass der Grad der
Unterstützung von HTML5 abhängig von der Einbindung in den Browser der
jeweiligen Browserhersteller ist. Jedes System mit einer aktuellen Version
eines Browsers kann damit die Funktionen von HTML5 verwenden und darstellen.
\newline\newline
Adobe Flash stellt das extreme Gegenteil zu HTML5 dar. Die Technologie ist
eine von Adobe entwickelte proprietäre Software und setzt vorraus, dass für
Browser ein Plug-In und für das System der Adobe Flash Player installiert sein
muss, um Flash-Inhalte darstellen zu können. Damit hat Adobe die Macht über die
Entscheidung der zu unterstützenden Plattformen. Sollte also Adobe sich dazu
entscheiden, eine bestimmte Plattform nicht mehr zu unterstützen, haben die
Nutzer dieser Plattform keine Möglichkeit mehr Flash-Inhalte zu konsumieren.
Apple hat sich aus diesem Grund dazu entschieden, auf ihren mobilen Engeräten
komplett auf die Unterstützung von Adobe Flash zu verzichten und ihre
Aufmerksamkeit auf offene Standard wie HTML5, CSS3 und JavaScript zu richten.
\newline\newline
Die gängigsten Desktop-Betriebssysteme (Windows, MacOS und Linux)
unterstützen sowohl HTML5, mittels diverser Browser, und Adobe Flash, mittels
Plug-In oder Player, und können damit Inhalte beider Technologien darstellen.
Bei mobilen Betriebssysteme sind immer vorinstallierte Browser vorhanden,
die einen Großteil der Features von HTML5 unterstützen. Für die Unterstützung
von Adobe Flash müssen zusätzliche Applikationen installiert werden.
\newline\newline
Um Flash-Inhalte auf Apples mobilen Endgeräte nutzen zu können, bietet Adobe
eine Softwarelösung, wodurch Flash Inhalte in HTML5 Inhalte umgewandelt werden.
Allerdings gilt es dabei zu beachten, dass die umgewandelten Inhalte nur
selten den originalen Inhalten entsprechen, aufgrund der fehlenden Funktionen
seitens HTML5 oder der hohen Komplexität der Inhalte. Adobe hat bereits
angekündigt, die Entwicklung des Adobe Flash Players für mobile Plattformen
einzustellen und sich auf Desktop Plattformen zu konzentrieren. Mit dieser
Entscheidung seitens Adobe in Zukunft der mobile Bereich sicherlich HTML5
gehören.

\section{Browserunterstützung}
Wie im letzten Abschnitt bereits beschrieben, ist die Darstellung von
HTML5-Inhalten abhängig von dem genutzten Browser. Unter den Browsern
gibt es eine breite Auswahl, die alle eine unterschiedlich fortgeschrittene
Unterstützung der Funktionen und Features von HTML5 bieten. Das führt
automatisch zu dem Problem, dass eine Webseite in unterschiedlichen
Browsern unterschiedlich dargestellt wird. Wie weit die Darstellung
dabei vom Original abweicht hängt von dem jeweiligen Browser ab.
Das stellt eines der größten Probleme im Bereich der Webentwicklung
mit den offenen Standards dar. Ein Entwickler muss sich darum bemühen,
dass eine möglichst große Anzahl an Nutzern die angebotene Seite/Applikation
optimal nutzen kann.
\newline\newline
An dieser Stelle soll daher ein Überblick darüber gegeben werden, welche
Browser welche Features und Funktionen von HTML5 unterstützen und der
Vergleich aufgestellt werden, ob diese mit Adobe Flash auch
erstellt/dargestellt werden können.
%caniuse grafik

\section{Leistungstests}
Bei der Auswahl einer Technologie ist neben den gebotenen Funktionen, dem damit
verbundenen Entwicklungs- und Lernaufwand und den Kosten auch die
Leistungsfähigkeit auf verschiedenen Betriebssystemen und Endgeräten
entscheidend. Speziell die Anzeigerate bei Videos und Animationen sowie die
Belastung des Prozessors spielen eine tragende Rolle, da die Anzeigerate
möglichst konstant und hoch und zeitgleich die Belastung des Prozessors
möglichst gering sein soll.
\newline\newline
Um diese Faktoren austesten zu können werden im Internet verschiedene
Tests angeboten, die eine bestimmte Situation simulieren und die
entsprechenden Ergebnisse liefern. Diese Tests, auch bekannt unter dem Namen
``Benchmarks'', decken dabei einen möglichst großen Testbereich ab. So reichen
diese von der Darstellung von Grafiken über das Berechnen von Partikeln
bis hin zu einem Belastungstest des Prozessors.
\newline\newline
Die Benchmark Ergebnisse können dabei lediglich als Richtwert angesehen
werden, da diese abhängig von dem beim Test verwendeten Betriebssystem und der
verwendeten Software (Browserhersteller und Version bzw. Version des Flash
Players) stark voneinander abweichen können.
\newline\newline
In den folgenden Abschnitten werden verschiedene Benchmarks verwendet, um
die Leistung von HTML5 und Adobe Flash möglichst genau zu vergleichen.
Bei den Benchmarks wird darauf geachtet, dass Funktionen verwendet werden,
die sowohl HTML5 als auch Adobe Flash zur Verfügung stellen damit für beide
Technologien die gleichen Bedingungen bestehen.

\subsubsection{Testumgebung}

Um sämtliche am Markt verfügbaren Browser zu testen, werden zwei
unterschiedliche Testsysteme verwendet. Die Testsysteme werden wie folgt
spezifiziert:
\begin{itemize}
  \item Gerät: Macbook Air (13 Zoll, Mitte 2012)
  \item Betriebssystem: OS X Version 10.9.1 (Mavericks)
  \item Prozessor: 2 GHz Intel Core i7
  \item Speicher: 8 GB 1600 MHz DDR3
  \item Grafikkarte: Intel HD Graphics 4000 1024 MB
  \item Festplatte: Flash-Speicher (SSD)
\end{itemize}
und
\begin{itemize}
  \item Gerät: HP Compaq nc6400 (13 Zoll)
  \item Betriebssystem: Windows 7 Professional, Service Pack 1
  \item Prozessor: 2 GHz Intel Core Duo 2 T7200
  \item Speicher: 2 GB
  \item Grafikkarte: ATI Mobility Radeon X1300 256 MB
  \item Festplatte: HDD
\end{itemize}
Die Auswahl der Browser wird wie folgt spezifiziert:
\begin{itemize}
  \item Firefox 26.0
  \item Chrome 32.0.1700.102
  \item Opera 19.0.1326.47
  \item Safari 7.0.1 (nur OS X)
  \item Internet Explorer 11 (nur Windows)
\end{itemize}
Die getestete Version vom Adobe Flash Player ist:
\begin{itemize}
  \item Adobe Flash Player 12.0.0.38
\end{itemize}

\subsection{GUIMark 2: The rise of HTML5}
% http://www.craftymind.com/guimark2/
GUIMark 2 ist ein kostenloser Benchmark entwickelt von Sean Christmann für
HTML5 und Adobe Flash. Christmann ist davon überzeugt, dass bei einer vollen
Ausschöpfung der Rendering Pipelines einfacher entschieden werden kann, welche
Technologie am Besten für interaktiven Inhalt im Web genutzt werden soll.
Laut Christmann tendieren Entwickler hauptsächlich dazu die Nutzbarkeit der
Technologien zu vergleichen, während in Wirklichkeit die meiste
Prozessorleistung für interne Rendering Schnittstellen genutzt wird.
\newline\newline
Der Benchmark ist in drei unterschiedliche Tests unterteilt: Vektor-, Bitmap-
und Textdarstellung. Die Testfälle sind dabei möglichst realen
Alltagsszenarien nachempfunden. Zusätzlich sind einige Testfälle vorhanden die
ebenfalls für mobile Endgeräte entwickelt wurden, um die Aussagen von Steve
Jobs, im Bezug auf die mobile Leistungsfähigkeit von HTML5 und Adobe Flash, zu
hinterfragen.
\newline\newline
Bei den folgenden Tests wird die Anzeigerate in Bilder pro Sekunde (fps,
Frames per Second) gemessen und miteinander verglichen.

\subsubsection{Vector Charting Test}
Dieser Benchmark wurde konzipiert um die Funktionen der Vektor APIs mittels
einem animierten Diagramm auszulasten. Das Diagramm wird dabei durch den
massiven Gebrauch von Linien und aufwendigen Transparenzfüllungen erstellt.
% Statistik
\begin{table}[H]
\begin{center}
\begin{tabular}{|p{3cm}| >{\centering\arraybackslash}p{3cm} | >{\centering\arraybackslash}p{3cm}|}

  \hline & HTML5 & Flash \\ \hline
  OSX 10.9.1 & & \\ \hline
  Firefox & 6,86 fps & 43,85 fps \\ \hline
  Chrome & 30,62 fps & 48,38 fps \\ \hline
  Opera & 29,86 fps & 46,98 fps \\ \hline
  Safari & \textbf{55.76 fps} & \textbf{49.96 fps} \\ \hline
  & i.D. 30,7 fps & i.D. 47,3 fps \\ \hline
  Windows 7 SP1 & & \\ \hline
  Firefox & 10,05 fps & 23,96 fps \\ \hline
  Chrome & 12,42 fps & \textbf{59,89 fps} \\ \hline
  Opera &  \textbf{13,52 fps} & 15,26 fps \\ \hline
  IE & 10,24 fps & 22,43 fps \\ \hline
  & i.D. 11,6 fps & i.D. 30,39 fps \\ \hline

\end{tabular}
\end{center}
\caption{Testergebnisse - Vector Charting Test}
\end{table}
Auf dem Macbook Air zeigt sich, dass Apples Browser Safari mit rund 55fps die
beste HTML5 Leistung bringt und die Leistung von Flash überall gleich konstant
bleibt.
% 1 Pixel Stroke Results
\begin{table}[H]
\begin{center}
\begin{tabular}{|p{3cm}| >{\centering\arraybackslash}p{3cm} |}

  \hline & HTML5 \\ \hline
  OSX 10.9.1 & \\ \hline
  Firefox & 44,91 fps \\ \hline
  Chrome & 54,43 fps \\ \hline
  Opera & 54,68 fps \\ \hline
  Safari & \textbf{56.92 fps} \\ \hline
  Windows 7 SP1 &  \\ \hline
  Firefox & 10.54 fps \\ \hline
  Chrome & 27.36 fps \\ \hline
  Opera & \textbf{30.38 fps} \\ \hline
  IE & 10.51 fps \\ \hline

\end{tabular}
\end{center}
\caption{Testergebnisse - Vector Charting Test 1 Pixel Stroke}
\end{table}
% Bild Test

\subsubsection{Bitmap Gaming Test}
Der Bitmap Test simuliert ein Spiel auf dem Prinzip des Towerdefense.
Dabei wird eine Vielzahl an Bitmaps pro Bild unterschiedlich animiert. Mit
jedem fertigen Bild muss dieses auch wieder gelöscht werden um alle neuen
Zustände im neuen Bild darstellen zu können. Um einen Tiefeneffekt zu erzeugen
werden Bitmaps auf unterschiedlichen Ebenen positioniert (Z Depth Ordering)
und um die Bitmaps zu skalieren wird Anti-alising verwendet.
% Statistik
\begin{table}[H]
\begin{center}
\begin{tabular}{|p{3cm}| >{\centering\arraybackslash}p{3cm} | >{\centering\arraybackslash}p{3cm}|}

  \hline & HTML5 & Flash \\ \hline
  OSX 10.9.1 & & \\ \hline
  Firefox & 20,9 fps & 36,95 fps \\ \hline
  Chrome & 55,61 fps & 35,02 fps \\ \hline
  Opera & 55,86 fps & 36,28 fps \\ \hline
  Safari & \textbf{57.03 fps} & \textbf{39.33 fps} \\ \hline
  & i.D. 47,35 fps & i.D. 36,9 fps \\ \hline
  Windows 7 SP1 & & \\ \hline
  Firefox & 7,34 fps & \textbf{20,09 fps} \\ \hline
  Chrome & 14,68 fps & 17,55 fps \\ \hline
  Opera & \textbf{18,71 fps} & 18,79 fps \\ \hline
  IE & 10,66 fps & 18,09 fps \\ \hline
  & i.D. 12,85 fps & i.D. 18,63 fps \\ \hline

\end{tabular}
\end{center}
\caption{Testergebnisse - Bitmap Gaming Test}
\end{table}
% Bild Test

\subsubsection{Text Column Test}
Bei diesem Test wird die Text Layout und Rendering Engine von HTML und Flash
ausgelastet. Es werden Custom Fonts, die mit CSS3 eingeführt wurden,
eingebunden und Multibyte Character Strings dargestellt. Anders als die
vorherigen Benchmarks stellt dieser Testfall kein reales Testszenario dar,
sondern dient der Analyse, wie viel Zeit benötigt wird, um eine textlastige
Webseite vollständig zu laden. Christmann nennt diesen Test den "`Eisberg"'-
Test, da ungefähr 80\% der Prozessorbelastung ausserhalb des Renderviews (des
sichtbaren Bereichs) entstehen. Dennoch soll der Test realistische Ergebnisse
liefern, da beim Scrollen berechnet werden muss, wieviele Zeilen des Textes nun
in bzw. aus dem Renderview bewegt werden müssen. Bei HTML Webseiten geschieht
dies immer, sobald eine Seite geöffnet wird und die Auflösung des genutzen
Monitors zu klein ist, um diese komplett anzeigen zu können.
% Statistik
\begin{table}[H]
\begin{center}
\begin{tabular}{|p{3cm}| >{\centering\arraybackslash}p{3cm} | >{\centering\arraybackslash}p{3cm}|}

  \hline & HTML5 & Flash \\ \hline
  OSX 10.9.1 & & \\ \hline
  Firefox & 28,09 fps & 32,3 fps \\ \hline
  Chrome & \textbf{31,85 fps} & \textbf{35,74 fps} \\ \hline
  Opera & 28,54 fps & 30,08 fps \\ \hline
  Safari & \textbf{31,85 fps} & 34,8 fps \\ \hline
  & i.D. 30,08 fps & i.D. 33,23 fps \\ \hline
  Windows 7 PS1 & & \\ \hline
  Firefox & \textbf{16,15 fps} & 11,53 fps \\ \hline
  Chrome & 15,29 fps & \textbf{18,74 fps} \\ \hline
  Opera & 11,39 fps & 11,38 fps \\ \hline
  IE & 10,6 fps & 11,08 fps \\ \hline
  & i.D. 13,36 fps & i.D. 13,18 fps \\ \hline

\end{tabular}
\end{center}
\caption{Testergebnisse - Text Column Test}
\end{table}
% Bild Test

\subsection{GUIMark 3: Mobile Showdown}

\subsection{The Man in Blue Animation Benchmark}

\subsection{Animation}

\subsection{Video}

\section{Besonderheiten}

% http://www.oesterreich.gv.at/site/5566/default.aspx#a4
% http://www.barrierefreies-webdesign.de/wcag2/
\subsection{Barrierefreiheit}
Barrierefreiheit im Internet bedeutet, Webseiten so zu gestalten, dass
jeder Mensch unabhängig von seinen körperlichen Einschränkungen die
Inhalte der Webseite nutzen kann. Dabei stellt besonders die Aufbereitung der
Inhalte für Menschen mit Sehbehinderungen eine Schwierigkeit dar.
Nichtsdestotrotz soll diese und auch andere Behinderungen bei der
Aufbereitung von Inhalten für das Internet berücksichtigt werden.
\newline\newline
Das W3C hat bereits im Jahre 1999 die erste Empfehlung für barrierefreies
Internet ausgesprochen. Für das Erstellen barrierefreier Webseiten werden,
seit 2008, Richtlinien in den Web Content Accessibility Guidlines (WCAG 2.0)
festgehalten und der Öffentlichkeit zugänglich gemacht. Die WCAG 2.0 bilden
nun einen flexibleren und testbareren, neuen Standard für barrierefreies
Webdesign. Dieser deckt Zugänglichkeitsanforderungen für alle
Arten von Web-Inhalten (Text, Bilder, Audio oder Video) sowie Web-
Applikationen ab und definiert technologieunabhängige Richtlinien und
Erfolgskriterien.
\newline\newline
Dabei sind die WCAG 2.0 pyramidenartig aufgebaut und umfassen vier Ebenen:
\begin{enumerate}
  \item 4 Prinzipien
  \item 12 Richtlinien
  \item 61 Erfolgskriterien
  \item unzählige Techniken
\end{enumerate}
Die ersten drei Ebenen stellen einen Maßstab und das Fundament der Richtlinien
dar. Im Gegensatz umfasst die vierte Ebene ergänzende Dokumente, die nicht
normativ sind und regelmäßig aktualisiert werden.
Im folgenden werden die Prinzipen und Richtlinien genauer betrachtet:
\begin{enumerate}
  \item Prinzip: Wahrnehmbar - Informationen und Bestandteile der
  Benutzerschnittstellen müssen den Benutzer so präsentiert werden, dass
  diese sie wahrnehmen können.
  \begin{itemize}
    \item Richtlinie 1.1 Textalternativen: Stellen Sie Textalternativen für
    alle Nicht-Text-Inhalte zur Verfügung, so dass in andere vom Benutzer
    benötigte Formen geändert werden können, wie zum Beispiel Großschrift,
    Braille, Symbole oder einfachere Sprache.
    \item Richtlinie 1.2 Zeitbasierte Medien: Stellen Sie Alternativen für
    zeitbasierte Medien zur Verfügung.
    \item Richtlinie 1.3 Anpassbar: Erstellen Sie Inhalte, die auf
    verschiedene Arten dargestellt werden können (z.B. einfacheres Layout),
    ohne dass Informationen oder Struktur verloren gehen.
    \item Richtlinie 1.4 Unterscheidbar: Machen Sie es Benutzern leichter,
    Inhalt zu sehen und zu hören einschließlich der Trennung von Vorder- und
    Hintergrund.
  \end{itemize}
  \item Prinzip: Bedienbar - Bestandteile der Benutzerschnittstelle und
  Navigation müssen bedienbar sein.
  \begin{itemize}
    \item Richtlinie 2.1 Per Tastatur zugänglich: Sorgen Sie dafür, dass alle
    Funktionalitäten per Tastatur zugänglich sind.
    \item Richtlinie 2.2 Ausreichend Zeit: Geben Sie den Benutzern ausreichend
    Zeit, Inhalte zu lesen und zu benutzen.
    \item Richtlinie 2.3 Anfälle: Gestalten Sie Inhalte nicht auf Arten, von
    denen bekannt ist, dass sie zu Anfällen führen.
    \item Richtlinie 2.4 Navigierbar: Stellen Sie Mittel zur Verfügung, um
    Benutzer dabei zu unterstützen zu navigieren, Inhalte zu finden und zu
    bestimmen, wo sie sich befinden.
  \end{itemize}
  \item Prinzip: Verständlich - Informationen und Bedienung der
  Benutzerschnittstellen müssen verständlich sein.
  \begin{itemize}
    \item Richtlinie 3.1 Lesbar: Machen Sie Inhalte lesbar und verständlich.
    \item Richtlinie 3.2 Vorhersehbar: Sorgen Sie dafür, dass Webseiten
    vorhersehbar aussehen und funktionieren.
    \item Richtlinie 3.3 Hilfestellung bei der Eingabe: Helfen Sie den
    Benutzern dabei, Fehler zu vermeiden und zu korrigieren.
  \end{itemize}
  \item Prinzip: Robust - Inhalte müssen robust genug sein, damit sie
  zuverlässig von einer großen Auswahl an Benutzeragenten einschließlich
  assistierender Techniken interpretiert werden können.
  \begin{itemize}
    \item Richtlinie 4.1 Kompatibel: Maximieren Sie die Kompatibilität mit
    aktuellen und zukünftigen Benutzeragenten, einschließlich assistierender
    Techniken.
  \end{itemize}
\end{enumerate}
Im Vergleich HTML5 und Adobe Flash ist es verständlich, dass die
Webtechnologie HTML einen Vorteil daraus zieht, dass die aufgezeigten
Richtlinien vor allem für das Erstellen von klassischen Webseiten gegeben
wurden. Allerdings bietet auch Adobe Flash verschiedene Möglichkeiten, um
diese Richtlinien erfüllen zu können. Diese sind jedoch mit deutlich höherem
Aufwand verbunden, als es bei HTML der Fall ist. Dieses Problem ist nicht nur
auf Adobe Flash beschränkt, sondern umfasst auch andere im Web verwendete
Technologien und Programmiersprachen, die bei der Umsetzung der Richtlinien für
Barrierefreiheit Schwierigkeiten bereiten. Das hat zur Folge, dass zum
Beispiel blinde Menschen große Schwierigkeiten bei der Nutzung von Nicht-HTML
Inhalten haben, da diese nicht von z.B. Screenreadern verarbeitet werden
können.
\newline\newline
Mit den aktuellen Versionen des Adobe Flash Players und der ActionScript
Version 3 bietet Adobe Entwicklern neue Möglichkeiten, um dem Nutzer
alternative Inhalte anzubieten. Grundsätzlich sind die Möglichkeiten dabei
ähnlich denen von HTML5. Damit wären barrierefreie Webseiten beziehungsweise
Anwendungen durchaus auch mit Adobe Flash umsetzbar. Trotzdem bleibt die
Umsetzung der Richtlinien des W3C ein schwieriges Unterfangen und ist deutlich
aufwändiger und somit auch teurer, sodass häufig darauf verzichtet wird.
\newline\newline
Es zeigt sich aus dem Vergleich HTML5 und Adobe Flash und bei der Betrachtung
der durch die W3C empfohlenen Richtlinien, dass Barrierefreiheit im Internet
nicht eine Frage der eingesetzten Technologie/n sondern der vorhandenen
Ressourcen (Zeit, Geld, Mitarbeiter, \dots) ist.

% Beispiele für Barrierefreien Code HTML5/Flash

\subsection{Suchmaschinenoptimierung}
Der Bereich der Suchmaschinenoptimierung hat sich in den letzten Jahren
enorm entwickelt. Kaum eine Webseite, unabhängig ob privat oder kommerziell
genutzt, kommt ohne sie aus. Das liegt vorwiegend daran, dass Suchmaschinen
wie Google, Yahoo oder Bing dem Nutzer die Suche nach interessanten und
relevanten Inhalte vereinfachen. Die Suchmaschinenoptimierung spielt dabei
insofern eine gewichtige Rolle, dass für viele Nutzer nur die Suchergebnisse
auf der ersten Seite relevant sind. Oft wird der Suchstring angepasst bevor
überhaupt Ergebnisse auf der zweiten oder dritten Ergebnisseite in Betracht
gezogen werden. Somit ist es für den Betreiber einer Webseite wichtig, dass
seine Webseite möglichst weit vorne im Suchmaschinenindex auftaucht.
\newline\newline
Inhalte die mit Adobe Flash erstellt wurden, haben einen großen
Nachteil: Crawler und Spider (die für die Indizierung und Bewertung der
Webseiten zuständig sind) können diese nicht oder nur sehr schwer lesen und
daher auch nicht indizieren. Die Indizierung ist allerdings notwendig, um
überhaupt auf einer Suchergebnisseite gelistet zu werden. Demzufolge ist es
nachvollziehbar, dass Webseiten und Inhalte, die auf Flash basieren, nicht in
das Ranking von Suchmaschinen mit aufgenommen werden. Die
Suchmaschinenbetreiber suchen immer wieder nach einer Lösung zu diesem
Problem. Die jüngsten Ergebnisse sind das Auslesen von Texten und Links aus
Flash-Daten oder externen Quellen wie XML-Dateien, die die Inhalte für die
Flash-Seiten beinhalten. Trotz allem ist die Indizierung bei Flash bei weitem
nicht so effektiv wie bei HTML. Bei einer multimedialen Flash-Anwendung ist
es zum Beispiel nicht möglich Bilder, Video- oder Audioinhalte und deren
Bedeutung zu extrahieren und die Zuordnung von Unterseiten zur Navigation der
Webseite ist nur äußerst umständlich über das sogenannte Deeplinking möglich.
Wobei dieses Problem auch bei anderen Webtechnologien, wie zum Beispiel AJAX,
mit dessen Hilfe Inhalte in die aktuell geöffnete Webseite nachgeladen werden
können, bestehen.
\newline\newline
Bei der Verwendung von Adobe Flash und der Suchmaschinenoptimierung bietet
Google einige Ratschläge, die helfen sollen: Suchmaschinen-Robotern sollen
stehts Zugriff auf die selben Inhalte wie die Nutzer haben. Das bedeutet,
dass keine zusätzlichen Inhalte aussschließlich für Crawler und Spider zur
Verfügung gestellt werden sollen. Die Angabe von Informationen die nur von
Suchmaschinen-Robotern gelesen werden können nennt man Cloaking. Grob
umschrieben gibt man den Robotern dabei vor etwas zu sein was man nicht ist.
Weitere Empfehlungen von Google sind, auf die Erstellung von vollständig
auf Flash basierenden Webseiten zu verzichten und stattdessen gezielt
einzelne Flash-Inhalte, zum Beispiel Videos, in eine HTML Seite zu
integrieren. Eine Alternative dazu wäre die gewünschten Inhalte sowohl mit
HTML5 als auch mit Adobe Flash zu erstellen, womit die Roboter die
HTML-Version indizieren können und die Nutzer die eventuell besser
aufbereitete Flash-Version besuchen können.
\newline\newline
Das ein Suchmaschinenbetreiber seinen Nutzern Ratschläge bei der Erstellung
von Flash-Inhalten bereitstellt, zeigt wie umständlich es sein kann, diese für
Suchmaschinen zu optimieren. Für eine möglichst erfolgreiche Optimierung
der Suchergebnisse sollte stets ein externes Datenformat wie HTML oder
XML eingesetzt werden. Damit ist es wesentlich einfacher, ein optimales
Suchmaschinenergebnis zu erhalten. Auch multimediale Inhalte können
durch einfache Mittel mit Bedeutung versehen werden, die wiederum für die
Crawler und Spider einfach verständlich sind. Unter anderem sind die weiteren
Vorteile bei der Nutzung von HTML die Möglichkeit, Hyperlinks unterschiedlich
stark zu bewerten, Meta-Attribute wie Stichwörter und/oder eine Beschreibung
der Seite anzugeben und damit eine differenziertere Indizierung zu erreichen
und somit ein deutlich besseres Ergebnis bei Suchmaschinen zu erreichen.
\newline\newline
Ist es für den Betreiber einer Webseite wichtig, im Ranking von Suchmaschinen
weit vorne zu liegen, führt fast kein Weg an HTML vorbei. Adobe Flash Inhalte
können, wie sich gezeigt hat, auch indiziert werden, allerdings sind deren
Relevanz für Suchmaschinen nie so hoch wie die von HTML Inhalten.
Mit der neuen Version von HTML stehen nun noch mehr Möglichkeiten zum Einbinden
von multimedialen und interaktiven Inhalten zur Verfügung, die wiederum
eine optimale Alternative für die Optimierung von Suchmaschinenergebnissen
genutzt werden können.

%Die neuen Möglichkeiten der neuen HTML-Version bieten dabei viele Vorteile bei
%der Suchmaschinenoptimierung wie zum Beispiel die Alternativen zur Einbindung
%von multimedialen und interaktiven Inhalten.

\subsection{Schnittstellen}
HTML5 und Adobe Flash sind mächtige Entwicklertools, welche die Realisierung
komplexer Anwendungen ermöglichen. Dennoch ist es mittlerweile im Internet
notwendig, dass Schnittstellen zu anderen Technologien entweder bereits
vorhanden oder zumindest realisierbar sind. Technologien, die oft in einem
Verbund verwendet werden, sind Datenbanken (SQL), Auszeichnungs- und
Austauschsprachen wie XML oder JSON oder Programmiersprachen wie PHP oder
JAVA.
\newline\newline
Da HTML5 eine Auszeichnungssprache und keine Programmiersprache ist, können
Inhalte und Daten aus anderen Technologien nur dargestellt werden. Mit
Hilfe von JavaScript lassen sich diese Daten allerdings auch verarbeiten.
Demnach ist im Fall von HTML5 nicht die Frage, ob die Technologie ausreichende
Schnittstellen bietet, sondern ob es genügend Möglichkeiten zur Integration
verschiedener andersartiger Technologien bietet und ob diese Technologien
auf JavaScript-Aufrufe reagieren können. Daher müssen die Daten der
eigesetzten Technologie so aufbereitet werden, dass sie mittels HTML
dargestellt werden können, zum Beispiel als Text, Bild oder sonstigen
multimedialen Inhalt.
\newline\newline
Adobe Flash bietet ebenfalls Möglichkeiten für den Austausch von Daten
zwischen verschiedenen Technologien an. Allerdings müssen, wie bei HTML5,
sämtliche Daten im Voraus aufbereitet werden um innerhalb von Flash
verarbeitet und dargestellt werden zu können. Mit der, in Flash integrierten
Programmiersprache ActionScript, lassen sich die Daten in vielfältigerer
Weise weiterverarbeiten, als es mit purem HTML5 möglich wäre. Dennoch kann
mit Flash nicht direkt auf andere Technologien zugegriffen werden.

%\section{Vergleich von HTML5 Canvas und Adobe Flash}
%Proprietär vs Standard
%\subsection{Verfügbarkeit}
%\subsection{Audio und Video}
%\subsection{Animation}
%\subsection{Spiele}
%\subsection{Werbung}
%\subsection{Web-Applikationen}

%\section{Mögliche Auswirkungen}
%\subsection{Der Wandel von dem allgegenwärtigen Flash zu den neuen Webstandards}
%\subsection{Die Zukunft von HTML5 und Flash}
%\subsection{Unternehmen und ihre Einstellung zu neuen Ideen}

%\chapter{Auswirkungen}
%\section{Browser}
%\subsection{Chrome und Safari (WebKit)}
%Chrome von Google und Safari von Apple bauen beide auf der HTML-Rendering-Engine
%WebKit zur Darstellung von Webinhalten auf. Google Chrome konnte zunehmend mit
%dem bis dato meist genutzen Webbrowser Internet Explorer konkurrieren, bis es
%nach den Angaben des globalen Statistiksunternehmens StatCounter erstmals im Mai
%2012, weltweit die Spitzenposition einnehmen konnte. Mit einem
%durchschnittlichen Anteil von 35\% für Google Chrome und 7\% für Apple Safari
%stellen diese die beliebtesten Webbrowser dar.
% http://gs.statcounter.com/#browser-ww-monthly-201204-201304-bar
%Im Vergleich decken beide Browser den größten Teil der HTML5 Unterstützung ab.
% Vergleiche HTML5 Unterstützung - http://html5readiness.com/

%Die WebKit Engine wurde auf der Grundlage des KDE-Projekts KHTML als Open Source
%Projekt von Apple entwickelt. Apple ist sehr daran interessiert, dass HTML5 so
%schnell wie möglich auf allen Geräten verfügbar ist, um den Benutzern ein
%uneingeschränktes Internet anbieten zu können. Auch die Mobile Safari Variante
%auf den iDivices basieren auf der WebKit-Engine.

%\subsection{Firefox (Gecko)}
%Der Open Source Webbrowser Firefox von Mozilla stellt im deutschsprachigen Raum
%den meist genutzten Webbrowser dar. Weltweit betrachtet befindet sich Firefox
%mit einem durchschnittlichen Anteil von 23\% hinter dem Internet Explorer von
%Microsoft und Google Chrome.
%Mozilla zeichnete sich schon früh als Unterstützer des HTML5 Standards ab. So
%war es schon mit sehr frühen Versionen von Firefox möglich, das Video-Element zu
%nutzen.

%\subsection{Internet Explorer (Trident)}
%Der wohl bekannteste Webbrowser Internet Explorer von Microsoft ist bereits in
%der zehnten Version veröffentlicht worden. Schon mit dem Internet Explorer 9 war
%Microsoft darauf bedacht die Geschwindigkeit ihres Browsers zu verbessern und es
%wurde darauf verzichtet, eigene Standards durchzusetzen.
%Erstmals orientierte man sich gänzlich an den Spezifikationen und unterstütze
%bereits einige HTML5 Features.
%So wurde das Audio/Video-Element unterstützt, wobei nur auf proprietäre Codecs
%wie H.264 und MP3/AAC zurückgegriffen wurde. Seit dem Internet Explorer 10 wird
%auch der freie Codec WebM unterstützt.

%Microsoft hatte sich aus seinem größten Manko einen Vorteil verschafft: Dadurch,
%das der Internet Explorer nur auf Windows Geräten läuft, kann er die ganze
%Plattform und Hardware nutzen. Im Gegensatz zur Browsersoftware der Konkurrenz,
%die auf den kleinsten gemeinsamen Nenner zurückgreifen müssen.
%So kann Internet Explorers JavaScript Engine Chakra, die Leistung von Multicore
%Prozessoren ausnutzen.

%Mit dem Internet Explorer 10 hat Microsoft einiges aus der Vergangenheit
%aufgeholt, allerdings sollten diese neuen Features nicht überbewertet werden,
%denn die Einsatzfähigkeit von HTML5 zeichnet sich weniger durch die neuen
%Browser, sonder viel mehr durch die Anzahl an veralteten Browser aus. Solange
%die alten Versionen des Internet Explorer im Gebrauch sind können viele
%Neuheiten von HTML5 nicht sinnvoll und ohne Fallbacklösung verwendet werden.
%# Implementierungen
\chapter{Implementierungen}

Um die in dieser Arbeit vorgestellten Technologien im praktischen Einsatz zu
veranschaulichen beschreibt dieses Kapitel die Implementierungen von drei
unterschiedlichen Applikationen. Diese Applikationen verwenden jeweils
verschiedene Features der HTML5 Spezifikation, wobei die Aufmerksamkeit
hauptsächlich an der Nutzung des Canvas Elements liegt, und Adobe Flash mit
ActionScript 3.0.

\section{Animation}

\subsection{Konzept}
\subsection{Umsetzung HTML5}
\subsection{Umsetzung Adobe Flash}

\section{Spiel}

\subsection{Konzept}
\subsection{Umsetzung HTML5}

  \subsubsection{Dokumente und Objekte}
  Bei der Umsetzung dieses Spiels werden HTML, CSS und JavaScript-Dateien
  benötigt. Zusätzlich werden Grafiken benötigt, die die Juwele im Spiel
  repräsentieren soll.

  1. Beginn
  Index.html - Benötigte CSS+JS Files festlegen.
  settings festlegen
  yepnope preload/loader

\subsection{Umsetzung Adobe Flash}

\section{Multimedia-Anwendung}

\subsection{Konzept}
\subsection{Umsetzung HTML5}
\subsection{Umsetzung Adobe Flash}


\section{HTML5 Frameworks for Animations, Games and more}
\begin{itemize}
  \item kineticjs (http://www.kineticjs.com) f
  \item fabricjs (http://fabricjs.com/) f
  \item paperjs (http://paperjs.org/) f
  \item createjs (http://www.createjs.com) f
  \item pixijs (http://www.pixijs.com/) f
  \item processingjs (http://processingjs.org/) f
  \item impactjs (http://impactjs.com/) €
  \item limejs (http://www.limejs.com/) f
  \item craftyjs (http://craftyjs.com/) f
  \item canvasengine (http://canvasengine.net/) f
  \item threejs (http://threejs.org/) f
  \item bhive (http://www.bhivecanvas.com/) f
  \item cakejs (https://code.google.com/p/cakejs/) f
  \item phaser.io (http://phaser.io/) f
  \item twojs(http://jonobr1.github.io/two.js/) f
\end{itemize}
\chapter{Zukunftsaussicht}

\section{Wird HTML5 Flash ersetzen?}

Steve Jobs deutet in seinem offenen Brief immer wieder an, dass
Flash veraltet und HTML5 die Zunkunft ist. Dies ist insofern falsch,
da HTML5 mit seiner bisherigen Browserunterstützung noch zu weit
entfernt ist, um mit dem Flash Player mitzuhalten. Unrecht hat
Steve Jobs mit der Aussage der fehlenden Offenheit allerdings nicht.
Flash soll eingesetzt werden, wenn das Ziel mit einem offenen
Standard nicht erreicht werden kann. Mit Flash kann qualitativ
hochwertige Leistung erbracht werden. Aufgrund seiner
Verbreitung ist es ein inoffizieller Standard.
\newline\newline
HTML5 steckt noch in den Kinderschuhen und hat es schwer, mit Flash
zu konkurrieren. Denn auch wenn Flash geschlossen ist, ist die Verbreitung
des Flashplayer immer noch höher als die Verbreitung von HTML5-fähigen
Browsern. Außerdem hat es ein Unternehmen wie Adobe wesentlich
einfacher, ein neues Feature in ihren Player einzuführen, als Konsortien wie
die W3C oder WHATWG, welche jedes Feature auf ihre Art demokratisch
einführen und dann auf Browsersupport hoffen.
\newline\newline
Im Videbereich kann sich HTML5 zwar schon präsentieren, doch auch auf
YouTube wird der Flash Player weiterhin eine wichtige Rolle spielen, da es
noch an wichtigen Funktionen in HTML5 Video fehlt, wie z.B. an einem
nativen Vollbild, einem Schutz vor Download zur Wahrung der Urheberrechte
oder der Kommunikation mit dem Nutzer. YouTube bietet dem Nutzer die
Möglichkeit, direkt über die Webcam ein Video hochzuladen, was ohne
Flash nicht möglich wäre.
\newline\newline
Entscheidend dafür, was demnächst in HTML5 und was weiterhin in Flash
umgesetzt wird, werden sicher die Interessen der Webentwickler sein.
Während der Anwender nämlich nur die Technik nutzt und es ihm
egak sein wird wie sein Video oder seine Webseite läuft, liegt es an dem
Entwickler, ob er überhaupt in relativ funktionsarmen Javascript statt
des sehr umfangreichen ActionScript 3 programmiert. Große Unternehmen
wie Yahoo, Facebook oder Google könnten ihre Vorteile wiederum
aus Webstandards ziehen, da sie selbst eigene Anforderungen an eine
neue Spezifikation einbringen können, ohne dass sie von einem
kommerziellen Anbieter wie Adobe abhängig sind.
\newline\newline
Die Frage, ob HTML5 den Flash Player in Zukunft ablösen wird, bleibt also offen.
Flash wird uns sicherlich noch eine Weile erhalten bleiben, da mit HTML5
in seiner jetztigen Form und Verbreitung noch nicht diesselbe Masse
erreicht werden kann. Sicher ist auch, dass keines der beiden Programme
das andere vollständig ersetzen kann, weshalb man auch nicht von
Konkurrenz sprechen sollte, sondern von gegenseitiger Ergänzung.
\newline\newline
So eignet sich Flash in Zukunft vo allem für Elemente, die eine hohe Performance
benötigen und auf eine umfrangreiche Scriptbibliothek (ActionScript 3)
zugreifen können, wie z.B. Spiele, Rich Internet Applications, komplexe
3D-Animationen, Audio/Videoplayer, Simulationen oder umfangreiche
Präsentationen mit hohem Audio- und Video-Anteil. Die HTML5/JavaScript
Funktionalitäten eigenen sich hingegen für interaktive Webseitenelemente
(Accordions, Tooltips, Dropdown-Menüs, Tabs, uvm), Formularvalidierung,
Chats oder einfache Präsentationen.

\section{HTML6}

Die Unterstützung von HTML5 ist noch nicht gewährleistet - hier und da wird
an allen Ecken noch gebastelt. Aber trotzdem hat im Januar der
Google-Mitarbeiter und das WHATWG-Mitglied Mark Pilgrim im WHATWG-Blog
bereits einen Einblick in HTML6 gewährt. Unter anderem wurde ein neues
Element mit dem Namen \em{<device>} vorgestellt, welches z.B.
Webcam-Konferenzen ohne den Umweg über Flash ermöglichen würde.
Aber ob diese Spezifikation auch wirklich "`HTML6"' heißen wird, ist unklar.
Weiter schreibt Pilgrim nämlich:

\begin{quote}
	The next version of HTML doesn't have a name yet. In fact,
	it may never have a name, because the working group is
	switching to an unversioned development model. Various
	parts of the specification will be at varying degrees of
	stability, as noted in each section. But if all goes according
	to plan, there will never be One Big Cutoff that is frozen in
	time and dubbed "`HTML6"'. HTML is an unbroken line
	stretching back almost two decades, and version numbers
	are a vestige of an older development model for standards
	that never really matched reality very well anyway.
	HTML5 is so last week. Let's talk about what's next.
\end{quote}

Pilgrim sagt also, dass die Vergabe von Versionsnummern veraltet ist und
besonders für den Veröffentlichungsprozess von HTML-Standards nicht funktioniert.
Ebenfalls interessant ist, dass der HTML5 WHATWG-Entwurf seit diesem
Eintrag seinen Namen immer wieder von "`WHATWG HTML (including HTML5)"'
zu "`HTML5 (including next generation addition still in development)"' ändert.
\newline\newline
All dies könnte darauf hindeuten, dass HTML5 wirklich die letzte versionierte
HTML-Spezifikation ist. Die Zukunft heißt also nicht "`HTML6"', sondern
einfach nur "`HTML"' oder "`HTML5 mit zusätzlichen Elementen"'. Das
würde die Entwicklung vorantreiben. Es muss nicht mehr ein ganzes
Paket an Neuerungen herausgebracht werden, auch einzelne Elemente
können eingeführt werden. Schließlich sind es die Browserhersteller, die
im Endeffekt HTML5 entwerfen. Dieses könnte eventuell die Existenz
der W3C in Zukunft überfküssig machen. Die beiden Spezifikationen sind
beinahe identisch und alles, was sich in der W3C-Version findet, steht
auch in der WHATWG-Version. Iam Hickson hat allerdings in einer E-Mail
kurz erwähnt, dass geplant ist, die WHATWG-Spezifikation in Zukunft ohne
engere Absprachen mit dem W3C zu erweitern. Das neue \em{<device>}
Element ist also nur ein Anfang.

\begin{quote}
	I've given up trying to keep a WHATWG copy of the
	HTML5 spec that matches what the W3C publish [...]
\end{quote}

Auch Buchautor Peter Kröner prophezeit, dass es eventuell nicht einmal zu
einer W3C-Version von HTML5 kommt.

\begin{quote}
	Die spannende Frage ist, ob es den "`One Big Cutoff"', sprich den
	Recommendation-Status für HTML5 zu einem Zeitpunkt, an dem
	es noch jemanden interessiert, zumindest für HTML5 geben wird.
	Ich glaube nicht.
\end{quote}

Dies begründet Kröner damit, dass die Ausarbeitung von HTML5 seit
dem Beginn der Entwicklung im Jahr 2004 bis heute (2013) neun
Jahre gedauert hat und inzwischen zum großen Teil implementiert
und benutzbar ist. Trotzdem wird aber von einem Recommendation-Status
gesprochen, der erst im Jahr 2022 angesetzt wird. Wenn im Jahr 2013
nun aber bereits von einem HTML6 \em{<device>} Element geredet wird,
könnte bereits im Jahr 2016 ein Standard namens HTML6 entwickelt und
großteilig in den Browsern implementiert sein und infolgedessen den
Stand von HTML5 im Jahr 2013 haben. Also ein Datum, was neun Jahre vor
dem Recommendation-Status von HTML5 liegt und im Jahr 2022 wahrscheinlich
niemanden mehr interessieren wird. Auch aus dieser Sicht scheint ein
versionsloses HTML-Modell mehr Sinn zu machen und ist laut
Kröner bereits auf dem Weg:

\begin{quote}
	Wenn die WHATWG obendrein plant, mit HTML6 ein neues,
	versionsloses Entwicklungsmodell einzuführen, betrifft dies
	bereits HTML5. Die Spezifikationen der WHATWG sind schon
	jetzt ein Mischmasch aus alten HTML-Features, Neuerungen,
	die sich auf einen breiten Konsens stützen hoch kontroversen
	Ideen wie HTML5-Microformats und dem \em{<device>}
	Element. Der unversionierte Ausbau ist also bereits im Gange."
\end{quote}
\include{tex/fazit}


% Literaturverzeichnis
% Das Literaturverzeichnis kann auch nach einem allfälligen Anhang positiioniert
% werden (siehe "`Leitfaden für Bachelor- und Diplomarbeiten"', Version 2.0,
% Abschnitt 2.9).

% Möglichkeit 1: Erzeugung des Literaturverzeichnisses mit BibTeX:
% Die Quellen sind in der Datei *.bib (hier Literatur.bib) einzugeben. Danach
% muss diese Vorlage einmal geTeXt werden, dann BibTeX angewendet werden und
% anschliessend nochmals zweimal geTeXt werden.
% Im Text erfolgt die Zitierung mit dem Anker-Schlüsselwort, z.B. \cite{kop05}.
\bibliographystyle{IEEEtran}
\nocite{*}
\bibliography{Literatur}


% Möglichkeit 2: Erzeugung eines Literaturverzeichnisses ohne BibTeX:
%\begin{thebibliography}{99}
%\bibitem[kop05]{kop05}
%H.~Kopka, {\em LaTeX, Band 1: Einführung}, Pearson Studium, München,
%3.~Auflage, 2005.
%\bibitem[knu98]{knu98}
%F.~Mittelbach, M.~Goossens, J.~Braams, D.~Carlisle, and Ch. Rowley, {\em The
% LaTeX Companion},
%Addison-Wesley, 2nd edition, 2004.
%\end{thebibliography}

% Abbildungsverzeichnis
\listoffigures
\addcontentsline{toc}{chapter}{Abbildungsverzeichnis} % fügt den Eintrag
% "Abbildungsverzeichnis" im Inhaltsverzeichnis hinzu
\newpage

% Tabellenverzeichnis
\listoftables
\addcontentsline{toc}{chapter}{Tabellenverzeichnis} % fügt den Eintrag
% "Tabellenverzeichnis" im Inhaltsverzeichnis hinzu
\newpage

% Abkürzungsverzeichnis
% Bei Verwendung der Dokumentklasse "scrartcl" ist der Befehlt
% \addchap{Abkürzungsverzeichnis} durch
% \addsec{Abkürzungsverzeichnis} zu ersetzen
\addchap{Abkürzungsverzeichnis}
\hspace{-17mm}\begin{tabular}{>{\raggedleft}p{0.2\linewidth} p{0.75\linewidth}
p{0.1\linewidth}}
%#
3G & 3rd Generation Mobile Telecommunications \\
%a
Ajax & Asynchronous JavaScript and XML \\
AS3 & ActionScript3 \\
API & Application Programming Interface (Programmierschnittstelle) \\
%b
%c
CSS & Cascading Style Sheet \\
CSS3 & Cascading Style Sheet Level 3 \\
%d
DOM & Document Object Model \\
%e
%f
%g
%h
HTML & HyperText Markup Language \\
HTML5 & HyperText Markup Language Version 5 \\
%i
%j
JS & JavaScript \\
JSON & JavaScript Object Notation \\
%k
%l
%m
%n
%o
%p
%q
%r
RIA & Rich Media Application \\
%s
SGML & Standard Generalized Markup Language \\
SQL & Structured Query Language \\
SVG & Scaleable Vector Graphics \\
%t
%u
UA & User Agent \\
URL & Uniform Resource Locator \\
%v
%w
W3C & World Wide Web Consortium \\

\end{tabular}


% Anhänge
\begin{appendix}
\chapter[Erster Anhang]{Überschrift des ersten Anhangs}

Text Text Text Text Text Text Text Text Text Text Text Text Text Text Text Text
Text Text Text Text Text Text Text Text ...


\chapter[Zweiter Anhang]{Überschrift des zweiten Anhangs}

Text Text Text Text Text Text Text Text Text Text Text Text Text Text Text Text
Text Text Text Text Text Text Text Text ...

\end{appendix}


\end{document}
